% macro conflict, first definition kept: \newcommand{\defeq}{\overset{\text{déf}}{=}}
% macro conflict, first definition kept: \newcommand{\defeq}{\overset{\text{déf}}{=}}
\documentclass[11pt]{article}
\usepackage[no-math]{fontspec}
\usepackage[english,french]{babel}
\usepackage[leqno]{amsmath}
\usepackage{amssymb}
\usepackage{amsthm}
\usepackage{mathrsfs}
\usepackage{tikz-cd}
\usepackage{enumitem}
\usepackage{marginfix}
\usepackage{needspace}
\AddToHook{env/dcomment/before}{\par\Needspace{6\baselineskip}}
\setlength{\emergencystretch}{1.5em}
\AddToHook{env/thebibliography/begin}{\selectlanguage{english}\raggedright}

\newcommand{\printedpage}[1]{\marginpar{\footnotesize p.~#1}}
\newcommand{\msp}[1]{\marginpar{\footnotesize [#1]}}

\newcommand{\cE}{\mathcal{E}}
\newcommand{\cC}{\mathcal{C}}
\newcommand{\cM}{\mathcal{M}}
\newcommand{\Ob}{\operatorname{Ob}}
\newcommand{\Fl}{\operatorname{Fl}}
\newcommand{\Hom}{\operatorname{Hom}}
\newcommand{\End}{\operatorname{End}}
\newcommand{\Ind}{\operatorname{Ind}}
\newcommand{\Ens}{\mathrm{Ens}}
\newcommand{\pr}{\mathrm{pr}}
\newcommand{\id}{\mathrm{id}}
\newcommand{\card}{\operatorname{card}}
\newcommand{\cardred}{\operatorname{card}\,\text{\upshape réd}}
\newcommand{\defeq}{\overset{\text{\upshape déf}}{=}}
\newcommand{\red}{\operatorname{\text{réd}}}
\newcommand{\circled}[1]{\tikz[baseline=(c.base)]\node[draw,circle,inner sep=1pt](c){#1};}
\newcommand{\ux}{\underline{x}}
\newcommand{\Img}{\operatorname{Im}}
\newcommand{\Ker}{\operatorname{Ker}}
\newcommand{\redop}{\mathop{\textnormal{réd}}\nolimits}
\newcommand{\into}{\lhook\joinrel\longrightarrow}
\newcommand{\cQ}{\mathcal{Q}}
\newcommand{\uS}{\underline{S}}
\newcommand{\uT}{\underline{T}}
\newcommand{\U}{\mathfrak{U}}
\newcommand{\cN}{\mathcal{N}}
\newcommand{\cF}{\mathcal{F}}
\newcommand{\cA}{\mathcal{A}}
\newcommand{\Max}{\operatorname{Max}}
\newcommand{\Ann}{\operatorname{Ann}}
\newcommand{\incl}{\lhook\joinrel\longrightarrow}
\newcommand{\limind}{\mathop{\underrightarrow{\mathrm{lim}}}}
\newcommand{\limproj}{\mathop{\underleftarrow{\mathrm{lim}}}}
\begin{document}
% ---- from der-ch18_p161-165.tex
\setcounter{footnote}{174}

%=====================================================================
\printedpage{159}
% head: A. Grothendieck : Dérivateurs, Ch. XVIII, Édition des Universités de Montpellier II et Paris VII

\msp{page 419}
est $\pi$-accessible, quand $x_0$ est dans $\cE_{0\,\pi}$. Or par hypothèse sur $\pi$, il est $\pi$-accessible si $x_0$ est dans $\cE_{0\,\pi_0}$. Mais si $x$ est dans $\cE_{0\,\pi}$, on aura % [sic] x -> x_0, the formula that follows concerns x_0
\[
x_0 \;=\; \varinjlim_{i\in I} x_{0\,i} \qquad
\left\{
\begin{array}{l}
I\ \ \text{grande}\ \ \text{devant}\ \ \pi_0,\\
\text{les}\ \ x_{0\,i}\ \ \text{dans}\ \ \cE_{0\,\pi_0},\\
\card I \leqslant \pi
\end{array}
\right.
\]
(car $\pi_0$ est utile pour $\cE_0$ par hypothèse, et $\pi$ satisfait $\pi^{\pi_0} = \pi \geqslant \cardred \cE_{\pi_0}$). Comme $x_0 \longrightarrow H(x_0,x_1) : \cE_0 \longrightarrow \Ens$ satisfait à $\mathrm{L}_{\pi_0}$, pour tout $x_1$ dans $\cE_1$, on aura donc % [sic] x_0 \longrightarrow H(x_0,x_1) -> x_0 \longmapsto H(x_0,x_1), elsewhere such functors are written with \longmapsto
\[
H(x_0) \;\simeq\; \varprojlim_{i\in I} H(x_{0\,i})
\]
(où pour $x_0$ dans $\cE_0$, $H(x)$ désigne le foncteur $x_1 \longmapsto H(x_0,x_1) : \cE_1 \longrightarrow (\Ens)$). Comme % [sic] H(x) -> H(x_0), the functor depends on x_0
les foncteurs $H(x_{0\,i})$ sont $\pi$-accessibles, et $\card I \leqslant \pi$, il s'ensuit que leur $\varprojlim\limits_{I}$ $H(x_0)$ est $\pi$-accessible (SGA 4 I 9.6 (ii)). Cela achève de prouver donc 2°).

3°). Pour prouver que $\Ind_\pi(\cE(\pi)) \longrightarrow \cE$ est essentiellement surjectif, il suffit de voir que tout $x \in \Ob\cE_0$ et tout $x \in \Ob\cE_1$
\msp{page 420}
est dans l'image essentielle. Or cela résulte de l'hypothèse que $\pi$ est utile pour $\cE_0$ et $\cE_1$, et de 1°) qui implique que $\cE_0 \longrightarrow \cE$ et $\cE_1 \longrightarrow \cE$ satisfont $\mathrm{L}_\pi$.

Cela achève de prouver la proposition 4.14.24, et de même le résultat suivant, plus précis que l'implication c) $\Longrightarrow$ a).

\textbf{Corollaire 4.14.25.} \textit{Soient $\cE$ une catégorie, $\cE_0$ une sous-catégorie ouverte, $\cE_1$ la sous-catégorie fermée complémentaire,}
\[
H \;=\; ((x_0,x_1) \longmapsto \Hom_{\cE}(x_0,x_1)) \;:\; \cE_0^\circ \times \cE_1 \;\longrightarrow\; \Ens
\]
\textit{le foncteur recollement.}
\begin{enumerate}
\item[\textit{1°)}] \textit{Soit $\pi$ un cardinal infini. Pour ce que $\cE$ et $\cE_0$ satisfassent à $L_\pi$, et le foncteur d'inclusion $\cE_0 \hookrightarrow \cE$ itou, il faut et il suffit que $\cE_0$ satisfasse à $L_\pi$, et que pour tout $x_1$ dans $\cE_1$, le foncteur $x_0 \longmapsto H(x_0,x_1) : \cE_0 \longrightarrow \Ens^\circ$ satisfasse à $L_\pi$. Alors $\cE_1$ et $\cE_1 \hookrightarrow \cE$ satisfont aussi à $L_\pi$,}
\msp{page 421}
\textit{et on a}
\[
\cE_\pi \;\subset\; \cE(\pi) \;\defeq\;
\begin{array}{l}
\text{\itshape sous-catégorie\ \ strictement\ \ pleine\ \ de\ \ $\cE$\ \ telle\ \ que}\\
\Ob\cE(\pi) = \Ob\cE_{0\,\pi} \cup \Ob\cE_{1\,\pi}.
\end{array}
\]
\item[\textit{2°)}] \textit{Supposons que la condition c) de 4.14.24 soit satisfaite, et soit $\pi_0$ un cardinal infini tel que $\pi_0$ soit utile pour $\cE_0$, et que les foncteurs $x_0 \longmapsto H(x_0,x_1) : \cE_0 \longrightarrow \Ens^\circ$ définis par les $x_1 \in \Ob\cE_1$ soient $\pi_0$-accessibles. Soit $\pi$ un cardinal tel que $\pi^{\pi_0} = \pi \geqslant \sup(\pi_0, \cardred\cE_{0\,\pi_0})$, et tel que pour tout $x_0 \in \Ob\cE_{\pi_0}$, le foncteur $x_1 \longmapsto H(x_0,x_1) : \cE_1 \longrightarrow \Ens$ soit $\pi$-accessible. Alors on a}
\[
\cE_\pi \;=\; \cE(\pi)\ .
\]
% foot: Version 17 décembre 2025

%=====================================================================
\printedpage{160}
% head: A. Grothendieck : Dérivateurs, Ch. XVIII, Édition des Universités de Montpellier II et Paris VII

\item[\textit{3°)}] \textit{Si $\pi_0$ et $\pi$ sont comme dans 2°), et si de plus $\pi$ est utile pour $\cE_1$, alors $\pi$ est utile pour $\cE$, donc (en vertu de 2°)}
\[
\cE \;\xleftarrow{\ \sim\ }\; \Ind_\pi(\cE(\pi))\ .
\]
\end{enumerate}
Nous pouvons maintenant énoncer la généralisation cherchée de 4.14.17 :

\msp{page 422}
\textbf{Théorème 4.14.26} (généralisation de 4.14.13). \textit{Soit $p : \cE \longrightarrow \cC$ un foncteur, avec $\cC$ petite. On suppose que $p$ satisfait à l'hypothèse H) de la page 409, qui équivaut (moyennant 4.14.24) à la suivante :}
\begin{enumerate}
\item[\textit{H$'$)}] \textit{$\forall\ s \in \Ob\cC$, la catégorie fibre $\cE_s$ est accessible, et pour tout $\lambda : s \longrightarrow t$ dans $\Fl\cC$, le foncteur}
\[
\left\{
\begin{array}{lcl}
H_\lambda : \cE_s^\circ \times \cE_t & \longrightarrow & \Ens\\[1ex]
\qquad (x_0,x_1) & \longmapsto & \Hom_\lambda(x_0,x_1)
\end{array}
\right.
\]
\textit{satisfait les deux conditions suivantes :}
\begin{enumerate}
\item[\textit{a)}] \textit{$\forall\ x_0 \in \Ob\cE_s^\circ$, $x_1 \longmapsto H_\lambda(x_0,x_1) : \cE_t \longrightarrow \Ens$ est accessible.}
\item[\textit{b)}] \textit{Il existe un cardinal infini $\pi_\lambda$, tel que pour tout $x_1 \in \Ob\cE_t$, le foncteur $x_0 \longmapsto H_\lambda(x_0,x_1) : \cE_s \longrightarrow (\Ens)^\circ$ soit $\pi_\lambda$-accessible.}
\end{enumerate}
\end{enumerate}
\textit{Ceci posé, la catégorie $\underline{\Gamma}(\cE/\cC)$ des sections de $\cE$ sur $\cC$ est accessible, et les foncteurs $\varepsilon_s : \underline{\Gamma}(\cE/\cC) \longrightarrow \cE_s$ sont accessibles. On peut préciser de façon similaire à 4.14.18 (énoncé en forme laissé au lecteur).}

\msp{page 423}
\textsc{Démonstration} de 4.14.26. Elle se fait exactement comme celle de 4.14.17 (p. 400 ff.), mais est conceptuellement plus simple. On suppose $\cC$ petite, et on introduit les catégories
\[
\Pi \;=\; \prod_{s\,\in\,\Ob\cC} \cE_s\ , \qquad \Pi' \;=\; \prod_{\lambda\,\in\,\Fl\cC} \cE_\lambda\ ,
\]
où
\[
\cE_\lambda \;=\; \cE \times_{\cC} (\mathbf{\Delta}_1, i_\lambda)\ ,
\]
et les deux foncteurs canoniques
\[
\Pi \;\underset{G}{\overset{F}{\rightrightarrows}}\; \Pi'
\]
de composantes
\[
\begin{aligned}
F_\lambda \;&=\; i_\lambda \circ \pr_{s(\lambda)}\\
G_\lambda \;&=\; j_\lambda \circ \pr_{b(\lambda)}\ ,
\end{aligned}
\]
où pour $\lambda : s \longrightarrow t$, on pose
\[
s(\lambda) \;=\; s\ , \qquad b(\lambda) \;=\; t\ ,
\]
et $i_\lambda$, $j_\lambda$ sont les inclusions canoniques
% In the print, the vertical arrow p_lambda and Delta_1 are enclosed together in a pair of large parentheses.
\[
\begin{tikzcd}
\cE_s \arrow[r,"i_\lambda"] \arrow[d,dash,"\wr"] & \cE_\lambda \arrow[dd,"p_\lambda"] & \cE_t\ . \arrow[l,"j_\lambda"'] \arrow[d,dash,"\wr"] \\
(\cE_\lambda)_0 & & (\cE_\lambda)_1 \\
& \mathbf{\Delta}_1 &
\end{tikzcd}
\]
% foot: Version 17 décembre 2025

%=====================================================================
\printedpage{161}
% head: A. Grothendieck : Dérivateurs, Ch. XVIII, Édition des Universités de Montpellier II et Paris VII

\msp{page 424}
On désigne par $\cM'$ la catégorie
\[
\cM' \;=\;
\begin{array}{l}
\text{catégorie des couples } (X,u)\text{, avec}\\
X \in \Ob\Pi,\ u : F(x) \longrightarrow G(x),
\end{array}
\] % [sic] F(x), G(x) -> F(X), G(X), the object is named X
qui est constructible par 4.14.22.1 est accessible (car $\Pi$ et $\pi'$ et les foncteurs $F$ et $G$ le sont, par l'hypothèse H) faite sur $p$). On peut l'interpréter comme la catégorie des sections % [sic] \pi' -> \Pi', the category introduced is \Pi'
de $\mathrm{Diag}(\cE)$ sur $\mathrm{Diag}(\cC)$. On a un foncteur strictement pleinement fidèle
\[
\cM \;\defeq\; \underline{\Gamma}(\cE/\cC) \;\hookrightarrow\; \cM'\ ,
\]
qui permet d'interpréter $\cM$ comme une sous-catégorie strictement pleine de $\cM'$. On voit alors comme dans le cas de 4.14.17 que cette sous-catégorie est accessible dans $\cM'$ -- c'est essentiellement sorital. Il s'ensuit que $\cM$ est accessible et que les foncteurs $\cM \longrightarrow \cE_s$ le sont aussi, comme composés de
\[
\cM \;\overset{\text{\upshape acc.}}{\hookrightarrow}\; \cM' \;\overset{\text{\upshape acc.}}{\hookrightarrow}\; \Pi \;\xrightarrow{\ \ \pr_s\ \ \text{\upshape acc.}\ \ }\; \cE_s\ ,
\]

\hfill q.e.d.

\msp{page 425}

\subsection*{4.15\quad Catégorie accessible sur une autre}

\textbf{Proposition 4.15.1.} \textit{Soit $f : \cE \longrightarrow \cC$ un foncteur avec $\cC$ petite. On suppose de plus satisfaite la condition H (p. 409), qu'on explicite ici sous la forme équivalente H$'$ :}
\begin{enumerate}
\item[\textit{H 1)}] \textit{Pour tout $s \in \Ob\cC$, $\cE_s$ est accessible.}
\item[\textit{H 2)}] \textit{Pour toute $\lambda \in \Fl\cC$, $\lambda : s \longrightarrow t$, existe un cardinal infini $\pi_\lambda$ tel que pour tout $x_1 \in \Ob\cE_t$, le foncteur}
\[
x_0 \;\longmapsto\; \Hom_\lambda(x_0,x_1) \;:\; \cE_s \;\longrightarrow\; (\Ens)^\circ
\]
\textit{est $\pi_\lambda$-accessible ($\cE_s$ satisfaisant de plus à $L_{\pi_\lambda}$).}
\item[\textit{H 3)}] \textit{Pour toute $\lambda \in \Fl\cC$, $\lambda : s \longrightarrow t$, et pour tout $x_0 \in \Ob\cE_s$, le foncteur}
\[
x_1 \;\longmapsto\; \Hom_\lambda(x_0,x_1) \;:\; \cE_t \;\longrightarrow\; (\Ens)
\]
\textit{est accessible.}
\end{enumerate}
\textit{Alors les foncteurs d'inclusion $i_s : \cE_s \longrightarrow \cE$ sont accessibles. De plus, si $\cC$ ``est sans projecteurs'' (i.e. si tout endomorphisme $p$ d'un objet $s$ de $\cC$ tel que $p^2 = p$, est l'identité), alors $\cE$ est accessible, et $f : \cE \longrightarrow \cC$ est accessible.}

Cela résulte du résultat plus précis :
% foot: Version 17 décembre 2025

%=====================================================================
\printedpage{162}
% head: A. Grothendieck : Dérivateurs, Ch. XVIII, Édition des Universités de Montpellier II et Paris VII

\msp{page 426}
\textbf{Corollaire 4.15.2} (\footnote{Pour a) et b) on n'utilise que H 1 et H 2, à l'exclusion de H 3.}). \textit{Les hypothèses sur $f : \cE \longrightarrow \cC$ sont celles de 4.15.1 ($\cC$ petite, et $f$ satisfait $H' =$ H 1) 2) 3)).}
\begin{enumerate}
\item[\textit{a)}] \textit{Soit $s \in \Ob\cC$ et soit $\pi_1$ un cardinal infini tel que pour toute $\lambda : s \longrightarrow t$ (flèche de $\cC$ de source $t$), on ait $\pi_1 \geqslant \pi_\lambda$ (où $\pi_\lambda$ est défini dans l'énoncé de H 2). Alors le foncteur $i_s : \cE_s \longrightarrow \cE$ est $\pi_1$-accessible.} % [sic] source t -> source s, lambda : s -> t has source s (cf. p. 163 "de source s")
\item[\textit{b)}] \textit{On suppose $\cC$ sans projecteurs. Soit $\pi_1$ un cardinal $\geqslant \sup(\pi_\lambda, \card\cC)$. Alors $\cE$ satisfait $L_{\pi_1}$, et $f : \cE \longrightarrow \cC$ est $\pi$-accessible.}
\item[\textit{c)}] \textit{Supposons encore que $\cC$ soit sans projecteurs (p.ex. une catégorie ``rigide'', p.ex. une catégorie préordonnée). Soient $\pi_0$, $\pi_1$ des cardinaux tels que :}
\begin{enumerate}
\item[\textcircled{\footnotesize $1^o$}] $\forall\ s \in \Ob\cC$, $\pi_0$ est utile pour $\cE_s$ \textit{(un tel $\pi_0$ existe grâce à H 1 et au fait que $\cC$, donc $\Ob\cC$ petits).}
\item[\textcircled{\footnotesize $2^o$}] $\pi_1 \geqslant \sup(\pi_0, \card\cC, \sup_{\lambda\in\Fl\cC}\pi_\lambda, \sup_{s\in\Ob\cC} \cardred \cE_{s\,\pi_0})$.
\item[\textcircled{\footnotesize $3^o$}] \textit{$\forall\lambda \in \Fl\cC$, et $x_0 \in \Ob\cE_{s,\pi_0}$, le foncteur $H_\lambda(x_0)$ de H 3, $x_1 \longmapsto \Hom_\lambda(x_0,x_1) : \cE_t \longrightarrow \Ens$ est $\pi_1$-accessible.}
\end{enumerate}
\textit{Soit $\pi$ un cardinal tel que $\pi^{\pi_0} = \pi \geqslant \pi_1$. Alors par b), $\cE$ satisfait $L_{\pi_1}$, donc $L_\pi$, et $f : \cE \longrightarrow \cC$ est $\pi_1$-accessible. De plus,}
\[
i_s(\cE_{s\,\pi}) \;\subset\; \cE_\pi\ ,
\]
\textit{ou encore}
\[
\cE(\pi) \;\subset\; \cE_\pi\ , \tag*{$(*)$}
\]
\textit{où $\cE(\pi)$ est la sous-catégorie pleine de $\cE$ telle que}
\[
\Ob\cE(\pi) \;=\; \bigcup_{s\in\Ob\cC} \Ob\,(\cE_{s\,\pi})\ .
\]
\textit{Comme $\cE_\pi$ satisfait $L_{\pi_1}$ par b), donc a fortiori $L_\pi$, on a donc par $(*)$ un foncteur} pleinement fidèle
\[
\Ind_\pi(\cE_{s\,\pi}) \;\xrightarrow{\ \sim\ }\; \cE\ ,
\]
\textit{et ce foncteur est une équivalence de catégories (donc $\cE$ est accessible).}
\item[\textit{d)}] \textit{En fait, sous les hypothèses de c) sur $\cC$, $\pi_0$, $\pi_1$, $\pi$, on a $\cE(\pi) = \cE_\pi$ (non seulement $\cE(\pi) \subset \cE_\pi$).}
\end{enumerate}

\msp{page 427}
\textsc{Démonstration de 4.15.1 et 4.15.2.} Il est clair que 4.15.1 est contenu dans 4.15.2 a) et c). Nous allons donc prouver 4.15.2.

a) Soit $(x_i)_{i\in I}$ un système inductif dans $\cE_s$, avec $I$ grand devant $\pi_0$, et $x$ sa limite inductive dans $\cE_s$ (qui existe, car $\cE_s$ satisfait à $\mathrm{L}_{\pi_\lambda}$, donc à $\mathrm{L}_{\pi_0}$). Il faut prouver que c'est une limite dans $\cE$, donc que pour tout $y \in \Ob\cE$, la flèche canonique
\[
\Hom_{\cE}(x,y) \;\longrightarrow\; \varprojlim_{I} \Hom_{\cE}(x_i,y) \tag*{$(*)$}
\]
% foot: Version 17 décembre 2025

%=====================================================================
\printedpage{163}
% head: A. Grothendieck : Dérivateurs, Ch. XVIII, Édition des Universités de Montpellier II et Paris VII
est un isomorphisme. Or $\Hom_{\cE}(x,y)$ et les $\Hom_{\cE}(x_i,y)$ sont des objets de $\Ens/\Hom_{\cC}(s,t)$, où $s = f(x)$, $t = f(y)$, et les flèches de transition sont des flèches dans $\Ens/\Hom_{\cC}(s,t)$. Cela montre que la bijectivité de $(*)$ revient à celle de
\[
\Hom_\lambda(x,y) \;\longrightarrow\; \varprojlim_{I} \Hom_\lambda(x_i,y) \tag*{$(*)$}
\]
pour tout $\lambda \in \Hom_{\cC}(x,y)$ \texttt{[plutôt} $\Hom_{\cC}(f(x),f(y))$\texttt{]}. Pour $\lambda$ fixé, et $y$ fixé dans $\Ob\cE_t$, la bijectivité dans $(*)$, pour tout système inductif grand devant $\pi_1$ dans $\cE_s$, de limite $x$, signifie que le foncteur
\[
x \;\longmapsto\; \Hom_\lambda(x,y) \;:\; \cE_s \;\longrightarrow\; (\Ens)^\circ
\]
est $\pi_0$-accessible. Or c'est bien le cas, par hypothèse sur $\pi_1$, quel que soit la flèche $\lambda : s \longrightarrow t$ de source $s$, et quel que soit $y \in \Ob\cE_t$.

\msp{page 428}
Soit $(x_i)_{i\in I}$ un système inductif grand devant $\pi_1$, et $(s_i)_{i\in I}$ l'image de ce système dans $\cC$. Nous prouverons plus bas le

% In the left margin, beside the statement of Lemme 4.15.3 and beside the following paragraph "Mais attention ...", the edition prints a wavy marginal sign (one for each).
\textbf{Lemme 4.15.3} (\footnote{Cf. démonstration page 452.},\footnote{sous réserve de la validité de 4.8.13 (Mais c'est à présent établi, sous une forme suffisante, dans 4.15.10).}). \textit{Soit $\cC$ une catégorie, $\pi$ un cardinal, tels que} $\card\cC \leqslant \pi$\textit{, et que $\cC$ soit sans projecteurs. Alors pour tout système inductif $(s_i)_{i\in I}$ grand devant $\pi$ dans $\cC$, existe une partie cofinale $J$ dans $I$ telle que le système induit $(s_i)_{i\in J}$ soit} constant, \textit{i.e. tel que les $s_i$ soient égaux à un $s \in \cC$, et que les morphismes de transition soient tous égaux à} $\id_s$.

Mais attention, on s'appuie pour prouver 4.15.3 sur le théorème 4.8.13 (page 212-216), que je n'ai toujours pas prouvé. Dans le cas où $\cC$ est rigide (i.e. $\End_{\cC}(s) = \{\id_s\}\ \forall\ s \in \Ob\cC$), la démonstration est cependant immédiate, cf. plus bas.

\msp{page 429}
Pour prouver que $x = \varinjlim\limits_{I} x_i$ existe dans $\cE$, et commute à $f$, on peut remplacer $I$ par une sous-catégorie cofinale. Cela nous ramène, grâce au lemme, au cas où $(s_i)_{i\in I}$ est constant, donc au cas où le système inductif provient d'un système inductif dans une catégorie fibre $\cE_s$. Mais par a), on sait que sa limite $x$ dans $\cE_s$ est une limite dans $\cE$, d'où l'existence de cette dernière, et la commutation à $f$, puisque $\varinjlim s_i = s$ ( $\varinjlim$ d'un système inductif filtrant constant).

c) Il faut prouver que si $x \in \Ob\cE_s$ est $\pi$-accessible dans $\cE_s$, il l'est dans $\cE$, i.e. que pour tout système inductif $(y_i)_{i\in I}$ grande devant $\pi$, la flèche canonique
\[
\varinjlim_{i} \Hom(x,y_i) \;\longrightarrow\; \Hom(x,y) \tag*{$(**)$}
\]
est bijective. Or comme par 2°) $\pi \geqslant \pi_1 \geqslant \sup(\card S, \sup_\lambda \pi_\lambda)$, il résulte encore du lemme % [sic] card S -> card \mathcal{C}, no S is defined; the condition 2° bounds card \mathcal{C}
que, quitte à remplacer $I$ par une sous-catégorie cofinale, on peut supposer que le système inductif $(y_i)_{i\in I}$ provient d'un système dans une catégorie fibre $\cE_t$, et que $y$ est sa $\varinjlim$ dans $\cE_t$. Mais alors les ensembles $\Hom(x,y)$ et $\Hom(x,y_i)$, et les morphismes de transition entre celles-ci, sont encore dans $\Ens/\Hom(s,t)$, et on est ramené
% foot: Version 17 décembre 2025

% ---- from der-ch18_p166-170.tex
\printedpage{164}
% head: A. Grothendieck : Dérivateurs, Ch. XVIII, Édition des Universités de Montpellier II et Paris VII

\msp{page 430}à prouver que pour tout $\lambda : s \longrightarrow t$, la flèche
\[
\varinjlim_{i} \Hom_{\lambda}(x, y_i) \ \longrightarrow \ \Hom_{\lambda}(x, y)
\]
est bijective. Pour $x$ et $\lambda$ fixés, et des systèmes inductifs variables dans $\cE_t$, grands devant $\pi$, cela signifie que le foncteur
\[
H_{\lambda}(x) \ = \ (y \longmapsto \Hom_{\lambda}(x, y)) \ : \ \cE_t \ \longrightarrow \ \Ens
\]
est $\pi$-accessible, et ceci quel que soit
\[
x \ \in \ \Ob \cE_{s\,\pi} \ .
\]
Or c'est le cas si $s \in \Ob \cE_{s\,\pi_0}$, en vertu de l'hypothèse $3^{\circ}$) (alors $H_{\lambda}(x_0)$ est $\pi_1$-accessible, % [sic] $s \in \Ob \cE_{s\,\pi_0}$ -> $x_0 \in \Ob \cE_{s\,\pi_0}$, the condition concerns an object of $\cE_{s\,\pi_0}$
et a fortiori $\pi$-accessible). Il faut passer de là au cas où $x \in \Ob \cE_{s\,\pi}$. Or comme $\pi_0$ est utile pour $\cE_s$ par $1^{\circ}$, on a
\[
x \ = \ \varinjlim_{j \in J} x_j \ , \qquad x_j \in \Ob \cE_{s\,\pi_0},\ J \text{ grande devant } \pi_0 \ ,
\]
et comme par $2^{\circ}$, $\pi^{\pi_0} = \pi \geqslant \pi_1 \geqslant \sup(\pi_0, \card \red \cE_{s\,\pi_0})$, on en conclut qu'on peut prendre $J$ de sorte que
\[
\card J \ \leqslant \ \pi \ .
\]
D'autre part, en vertu de $\pi \geqslant \pi_1 \ \geqslant \ \pi_{\lambda}$, le foncteur
\[
\left\{
\begin{array}{lcl}
x_0 & \longmapsto & H_{\lambda}(x_0) \\[1ex]
\cE_s & \longrightarrow & (\underline{\Hom}(\cE_s, \Ens))^{\circ}
\end{array}
\right.
\]
est $\pi_0$-accessible, donc on a
\[
H_{\lambda}(x_0) \ \simeq \ \varprojlim_{j \in J} H_{\lambda}(x_j) \ .
\]

\msp{page 431}Comme $x_j \ \in \ \Ob \cE_{s\,\pi_0}$, $H(x_j)$ est un foncteur $\pi_1$-accessible, et comme $\card J \leqslant \pi$, la $\varprojlim\limits_{j \in J} H_{\lambda}(x_j)$ est aussi $\pi$-accessible (cf. SGA 4 I 9.8).

Ainsi on a prouvé $\cE_{s\,\pi} \subset \cE(\pi)$ pour tout $s \in \Ob \cC$, i.e. l'inclusion
\[
\cE(\pi) \ \subset \ \cE_{\pi} \ ,
\]
d'où le foncteur pleinement fidèle
\[
\Ind_{\pi}(\cE(\pi)) \longrightarrow \cE \ , \tag*{$(*)$}
\]
comme $\cE$ satisfait à $\mathrm{L}_{\pi}$. Il reste à prouver que c'est essentiellement surjectif. Mais si $x \in \Ob \cE$, $s = f(x)$, on a $x \in \Ob \cE_s$, et comme $\pi_0$ est utile pour $\cE_s$ par $1^{\circ}$), et $\pi^{\pi_0} = \pi > \pi_0$,
$\pi$ est également utile pour $\pi$, donc on a \textit{dans} $\cE_s$ % [sic] utile pour $\pi$ -> utile pour $\cE_s$, usefulness is a property relative to the category $\cE_s$
\[
x \ = \ \varinjlim_{I} x_i \ , \qquad \text{les } x_i \in \cE_{s\,\pi} \subset \cE(\pi),\ I \text{ grande devant } \pi \ ,
\]
\printedpage{165}
% head: A. Grothendieck : Dérivateurs, Ch. XVIII, Édition des Universités de Montpellier II et Paris VII
et d'après a) (comme $\pi \geqslant \pi_1 \geqslant \sup_{\lambda \in \Fl \cC} \pi_{\lambda}$) $x$ est aussi limite des $x_i$ dans $\cE$, donc il est dans l'image essentielle de $(*)$.

d) Reste à prouver que l'on a $\cE_{\pi} \ \subset \ \cE(\pi)$. Mais on sait (par 4.8.~\dots), comme $\cE \ \simeq \ \Ind_{\pi}(\cE(\pi))$, que $\cE_{\pi} \ = \ \cE(\pi)^{\natural}$, i.e. que pour $x \ \in \ \Ob \cE_{\pi}$, $\exists \ y \ \in \ \cE(\pi)$ et un projecteur $p$ dans $x$, tel que $x$ soit isomorphe à l'image de $p$ dans $\cE$. Or $f(p)$ est un projecteur dans $f(y) = s$, donc par hypothèse sur $\cC$, $f(p) = \id_s$, i.e. $p$ est un projecteur dans $\cE_s$. Comme donc $x$ est dans $\cE_s$, et

\msp{page 432}même dans $\cE_{s\,\pi}$, puisque $\cE_{s\,\pi}$ est stable dans $\cE_s$ par facteurs directs. Cela prouve que $x \in \Ob \cE(\pi)$, q.e.d.

On peut présenter 4.15.1 sous une forme plus sympathique.

\textbf{Corollaire 4.15.4.} \textit{Soit $f : \cE \longrightarrow \cC$ un foncteur de $\mathfrak{U}$-catégories. Les conditions suivantes sont équivalentes :}
\begin{enumerate}
\item[\textit{a)}] \textit{Pour tout foncteur $\cC' \longrightarrow \cC$, avec $\cC'$ petite et sans projecteurs, $\cE' \overset{\text{déf}}{=} \cE \times_{\cC} \cC'$ est accessible, et $f' : \cE' \longrightarrow \cC'$ est accessible.}
\item[\textit{b)}] \textit{Comme a), mais avec $\cC' = \mathbf{\Delta}_1$, en d'autres termes, c'est la condition $H$ p.~409 : pour toute $\lambda \in \Fl \cC$, $\cE_{\lambda} = \cE \times_{\cC} (\mathbf{\Delta}_1, i_{\lambda})$ est accessible, et $\cE_{\lambda} \longrightarrow \mathbf{\Delta}_1$ est accessible.}
\item[\textit{c)}] \textit{La condition $H' = H\ \mathit{1} + H\ \mathit{2} + H\ \mathit{3}$, explicitée dans l'énoncé de 4.15.1.}
\end{enumerate}
\textsc{Démonstration}. Elle se fait circulaire,
\[
\mathrm{a} \ \Longrightarrow \ \mathrm{b} \ \Longrightarrow \ \mathrm{c} \ \Longrightarrow \ \mathrm{a} \ ,
\]
l'implication a $\Longrightarrow$ b étant tautologique, b $\Longrightarrow$ c l'étant presque (cf. 4.14.24, a $\Longrightarrow$ b $\Longrightarrow$ c). En fait, les conditions b et c sont des conditions sur les catégories $\cE_s \xrightarrow{f_{\lambda}} \mathbf{\Delta}_1$ pour $\lambda \in \Fl \cC$, donc stables par changement de base. Pour prouver

\msp{page 433}c $\Longrightarrow$ a, on peut donc supposer que $\cC' = \cC$, $\cC' \longrightarrow \cC$ le foncteur identique, et prouver que si c) est satisfaite et si $\cC$ est petite et sans projecteurs, alors $\cE$ est accessible et $f$ accessible. Or c'est ce qu'affirme 4.15.1.

\textbf{Remarque 4.15.5.} Si $f : \cE \longrightarrow \cC$ est un foncteur pleinement fidèle, il est immédiat que la condition H est satisfaite (car $\cE_{\lambda} \longrightarrow \mathbf{\Delta}_1$ sera pleinement fidèle, donc $\cE_{\lambda}$ est équivalente
à $e$ où à $\mathbf{\Delta}_1$, donc accessible, et tout foncteur de $\cE_{\lambda}$ dans une catégorie quelconque est % [sic] où -> ou, disjunction "e ou Δ1"
accessible \dots). Mais si on suppose $\cC$ petite \textit{sans} supposer $\cC$ sans projecteurs, il ne s'ensuit pas que $\cE$ soit accessible, ou ce qui revient au même (étant essentiellement petite), que $\cE$ soit karoubienne. (L'exemple le plus trivial est celui où $\cE = \cC$, $f = \id_{\cC}$, $\cE$ n'est karoubienne que si $\cC$ l'est !) Mais même en supposant $\cC$ une petite catégorie karoubienne, il ne s'ensuit pas forcément que $\cE$ l'est. En effet, \textit{toute} petite catégorie $\cE$ (karoubienne ou non) se plonge comme sous-catégorie pleine dans son enveloppe karoubienne $\cC$ -- il suffit donc de prendre $\cE$ non karoubienne. L'exemple

\msp{page 434}le plus économique est
\[
\Ob \cE \ = \ \{0\} \ , \qquad \Fl(\cE) \ = \ \{\id_0, p\} \ , \qquad \text{avec } p^2 = p \neq \id_0 \ ,
\]
\printedpage{166}
% head: A. Grothendieck : Dérivateurs, Ch. XVIII, Édition des Universités de Montpellier II et Paris VII
l'enveloppe de Karoubi de $\cE$ est donnée par
\[
P \ = \ \left\{
\begin{tikzcd}
1 \arrow[r, shift right, "\sigma"'] & 0 \arrow[l, shift right, "\rho"'] \arrow[loop right, "p"]
\end{tikzcd}
\qquad
\begin{aligned}
\rho\sigma \ &= \ \id_1 \\
\sigma\rho \ &= \ p \ \neq \ \id_0
\end{aligned}
\quad \Longrightarrow \ p^2 = p
\right\} .
\]
Cela soulève la question de trouver un énoncé similaire à 4.15.1, donnant des conditions serviables simples, plus fortes que $\mathrm{H}' = \mathrm{H}$ (1+2+3), sur $f : \cE \longrightarrow \cC$, pour que $\cE$ soit accessible et $f$ accessible, en supposant $\cC$ \textit{karoubienne}. Ou encore, des conditions nécessaires et
suffisants serviables (plus fortes que 4.15.4) sur $f : \cE \longrightarrow \cC$, pour que toute \textit{petite} catégorie % [sic] suffisants -> suffisantes, agreement with "conditions"
\textit{karoubienne} $\cC'$, et tout foncteur $\cC' \longrightarrow \cC$, la catégorie $\cE' = \cE \times_{\cC} \cC'$ soit accessible, et le foncteur $\cE' \longrightarrow \cC'$ accessible. Une idée naturelle, c'est qu'il suffirait que cette conclusion soit vraie quand $\cC'$ est une des deux catégories $\mathbf{\Delta}_1$, $P$ (où $P$ est la catégorie à deux objets, à trois flèches non identiques, définie plus haut). Mais je doute qu'il en soit ainsi (\footnote[178]{Ça devrait marcher par contre si on prend les catégories $0 \longrightarrow 1$,
\[
\begin{array}{c}
2 \longrightarrow \underbrace{\begin{tikzcd}[ampersand replacement=\&]
1 \arrow[r, shift right, "\sigma"'] \& 0 \arrow[l, shift right, "\rho"'] \arrow[loop right, "p"]
\end{tikzcd}}_{P}
\quad \text{et} \quad
2 \longleftarrow \underbrace{\begin{tikzcd}[ampersand replacement=\&]
1 \arrow[r, shift right, "\sigma"'] \& 0 \arrow[l, shift right, "\rho"'] \arrow[loop right, "p"]
\end{tikzcd}}_{P}
\qquad (\text{\footnotemark[179]}) \ , \\
\nwarrow \ \begin{array}{c} \text{3 objets et} \\ \text{5 flèches non} \\ \text{identiques} \end{array} \ \nearrow
\end{array}
\]
grâce à 4.15.7, page 440.}\footnotetext[179]{ces catégories représentent respectivement les foncteurs $\mathrm{Cat} \longrightarrow \Ens$ associant à une catégorie $\cC$ les diagrammes dans $\cC$ formés de $(x, y, p, u)$, avec $p$ projecteur dans $p$ \texttt{[plutôt dans $\cC$]}, et $u : y \longrightarrow x^{\natural}$ (i.e. $u : y \longrightarrow x$ avec $u = pu$), resp. $u : x^{\natural} \longrightarrow y$ (i.e. $u : x \longrightarrow y$ avec $u = up$).}).

On va quand même essayer un énoncé un peu très particulier dans ce sens :

% Verified on PDF leaf8: a non-textual wavy marginal mark beside items a), b); no mathematical token omitted.
\msp{page 435}\textbf{Proposition 4.15.6} (\footnote[180]{\textit{Faux ?}}). \textit{Soit $f : \cE \longrightarrow \cC$ un foncteur de $\mathfrak{U}$-catégories. On suppose :}
\begin{enumerate}
\item[\textit{a)}] \textit{$f$ est fibrant ou cofibrant} (\footnote[181]{Le cas fibrant n'est pas prouvé, peut-être \textit{faux} -- et de même pour le cas cofibrant !}).
\item[\textit{b)}] \textit{Les catégories fibres $\cE_s$ ($s \in \Ob \cC$) sont accessibles, et pour toute flèche $\lambda : s \longrightarrow t$ dans $\cC$, le foncteur changement de base $\lambda^* : \cE_t \longrightarrow \cE_s$ (resp. cochangement de base $\lambda_* : \cE_s \longrightarrow \cE_t$) est accessible.}
\item[\textit{c)}] \textit{$\cC$ est essentiellement petite et karoubienne.}
\end{enumerate}
\textit{Alors $\cE$ est accessible, et $f$ est accessible.}

\textsc{Démonstration}. On a déjà vu que les conditions a) et b) impliquent $\mathrm{H}'$, i.e. H (1+2+3). D'ailleurs, si $\cC'$ est une catégorie et $\cC' \xrightarrow{\ \sim\ } \cC$ une équivalence, donc alors $\cE' \overset{\text{déf}}{=} \cE \times_{\cC} \cC' \longrightarrow \cE$ est une équivalence, donc la conclusion voulue pour $\cE$, $f$ équivaut à la même conclusion pour $\cE'$, $f' : \cE' \longrightarrow \cC'$. Donc on peut supposer que $\cC$ est petite (et même réduite, si on y tient). En vertu de 4.15.2 a), on voit donc que les foncteurs d'inclusion $\cE_s \longrightarrow \cE$ sont accessibles (et même $\pi_1$-accessibles) dès que $\pi_1 \geqslant \sup_{\lambda \in \Fl \cC} \pi_{\lambda}$. On voit d'ailleurs

\msp{page 436}par AQT \texttt{[âne qui trotte]} qu'on peut trouver des cardinaux $\pi_0$, $\pi_1$ ayant les propriétés suivantes :

\printedpage{167}
% head: A. Grothendieck : Dérivateurs, Ch. XVIII, Édition des Universités de Montpellier II et Paris VII
\begin{enumerate}
\item[\circled{1\textsuperscript{o}}] $\pi_0$ est utile pour $\cE_s$, $\forall \ s \in \Ob \cC$.
\item[$2^{\circ}$)] $\pi_1 \geqslant \sup(\pi_0, \card \cC, \sup_{\lambda \in \Fl \cC} \pi_{\lambda}, \sup_{s \in \Ob \cC} \card \red \cE_s)$.
\item[\circled{3\textsuperscript{o}}] $\forall \ (\lambda : s \longrightarrow t) \in \Fl \cC$ et $x_0 \in \Ob \cE_{s\,\pi_0}$, le foncteur
\[
H_{\lambda}(x_0) \ : \ (x_1 \longmapsto \Hom_{\lambda}(x_0, x_1)) \ : \ \cE_t \ \longrightarrow \ (\Ens)
\]
est $\pi_1$-accessible. (\textbf{N.B.} Si $f$ est cofibrant, cela signifie que $\forall \ \lambda \in \Fl \cC$, $\lambda : s \longrightarrow t$,
\fbox{$\lambda_*(\cE_{s\,\pi_0}) \ \subset \ \cE_{s\,\pi_1}$}. Si $f$ est fibrant, il suffit d'exiger que $\forall \ \lambda \in \Fl \cC$, $\lambda^* : \cE_t \longrightarrow \cE_s$ % [sic] $\cE_{s\,\pi_1}$ -> $\cE_{t\,\pi_1}$, $\lambda_*$ goes to $\cE_t$, cf. the later display "(contenu dans 3°)"
soit $\pi_1$-accessible.) Mais dans ce cas, on veut aussi \fbox{$\lambda^*(\cE_{s\,\pi_0}) \ \subset \ \cE_{s\,\pi_1}$}. % [sic] $\lambda^*(\cE_{s\,\pi_0})$ -> $\lambda^*(\cE_{t\,\pi_0})$, $\lambda^*$ is defined on $\cE_t$, cf. the later display "(contenu dans 3°)"
\item[\circled{4\textsuperscript{o}}] Les foncteurs $\lambda^*$ (cas $f$ fibrant) resp. $\lambda_*$ (cas $f$ cofibrant) sont $\pi_1$-accessibles.
\end{enumerate}
Soit alors $\pi$ un cardinal tel que
\[
\pi \ = \ \pi^{\pi_0} \ \geqslant \ \pi_1 \ , \tag*{$(*)$}
\]
et soit $\cE(\pi)$ la sous-catégorie (strictement) pleine de $\cE$ définie par
\[
\Ob \cE(\pi) \ = \ \bigcup_{s \, \in \, \Ob \cC} \Ob \cE_{s\,\pi} \ .
\]
Grâce à $(*)$, et à $1^{\circ}$, $4^{\circ}$ et à
\[
\left\{
\begin{array}{l}
\pi_1 \ \geqslant \ \sup_s(\card \red \cE_{s\,\pi_0}) \qquad \text{(contenu dans $2^{\circ}$)} \\[1ex]
\lambda^*(\cE_{t\,\pi_0}) \ \subset \ \cE_{s\,\pi_1} \qquad \text{resp.} \qquad \lambda_*(\cE_{s\,\pi_0}) \ \subset \ \cE_{t\,\pi_1} \qquad \text{(contenu dans $3^{\circ}$)} \ ,
\end{array}
\right.
\]
on voit que $\cE(\pi)$ est une sous-catégorie

\msp{page 437}fibrée, resp. une sous-catégorie cofibrée, sur $\cC$ (le foncteur $\cE(\pi) \longrightarrow \cE$ étant donc cartésien resp. cocartésien). À $\cC$-équivalence près, $\cE$ se reconstitue alors à partir de $\cE(\pi)$, en passant au $\Ind_{\pi}(\cE_{s\,\pi})$, qui forment un système fibré resp. cofibré sur $\cC$ par les foncteurs $\Ind_{\pi}(\lambda^*_{\pi})$ resp. $\Ind_{\pi}(\lambda^{\pi}_*)$ \dots

Prouvons que $\cE$ satisfait à $\mathrm{L}_{\pi}$ et que $f : \cE \longrightarrow \cC$ est $\pi$-accessible (\footnote[182]{Ici on utilise seulement $\pi \geqslant \card \red \cC$.}). Comme $\cC$ est karoubienne, et $\pi \geqslant \pi_1 \geqslant \card \cC$, $\cC$ satisfait à $\mathrm{L}_{\pi}$ et les ind-objets grands devant $\pi$ de $\cC$ sont essentiellement constants (4.8.13). Soit alors $I$ grande devant $\pi$, et $I \longrightarrow \cE$, $i \longmapsto x_i$, un foncteur, il faut construire une $\varinjlim$ dans $\cE$, en utilisant la $\varinjlim$ déjà connue $s$ des $s_i = f(x_i)$. On va prendre un $x$ dans $\cE_s$. Prenons d'abord le cas où $f$ est cofibrant. On voit que les $\beta_i(x_i)$ forment un système inductif dans $\cE_s$, puisque pour $i = j$, on a
\[
\beta_{i\,*}(x_j) \ = \ \beta_{j\,*} \underbrace{\beta_{ji\,*}(x_i)}_{\longrightarrow \, x_j} \ \longrightarrow \ \beta_{j\,*}(x_j) \ . % [sic] $\beta_{i\,*}(x_j)$ -> $\beta_{i\,*}(x_i)$, the object transported from $s_i$ is $x_i$
\]
\[
\begin{tikzcd}[row sep=small]
x_i \arrow[rr, "\alpha_i"] & & x \, ? \\
s_i \arrow[rr, "\beta_i"] \arrow[dr, "\beta_{ji}"'] & & s \\
& s_j \arrow[ur, "\beta_j"'] &
\end{tikzcd}
\]
On pose

\printedpage{168}
% head: A. Grothendieck : Dérivateurs, Ch. XVIII, Édition des Universités de Montpellier II et Paris VII
\msp{page 438}
\[
x \ = \ \varinjlim_{I} \beta_{i\,*}(x_i) \ , \qquad \varinjlim \ \textit{dans } \cE_s \ .
\]
On sait que c'est une $\varinjlim$ dans $\cE$, donc que pour tout $y$ dans $\cE$, on a
\[
\Hom(x, y) \xrightarrow{\ \sim\ } \varprojlim_{I} \Hom(\beta_{i\,*}(x_i), y) \ ,
\]
ou ce qui revient au même, pour toute flèche $\lambda : s \longrightarrow t$, et $y$ dans $\Ob \cE_t$, on a
\[
\underbrace{\Hom_{\lambda}(x, y)}_{= \, \Hom_{\cE_t}(\lambda_* x, y)} \ \simeq \ \varprojlim \Hom_{\cE_t}(\underbrace{\lambda_* \beta_{i\,*}(x_i)}_{(\lambda_i \beta_i)_*(x_i)}, y) \ .
\]
Cela résulte aussi du fait que $\lambda_*$ est $\pi$-accessible (étant $\pi_1$-accessible) \dots{} On vérifie facilement que c'est une $\varinjlim$ des $x_i$ dans $\cE$. Et bien sûr, cette limite commute avec $f$ par construction.

Mais dans le cas où $f$ est fibrant, je ne vois comment faire pour construire une limite dans $\cE_s$. Il est vrai qu'on peut supposer que tous les $s_i$ sont égaux à $s$, mais les transitions $\beta_{ji}$ ne sont pas l'identité, ni même des isomorphismes. Il faudrait prendre $\cC = P$ pour essayer de comprendre -- je ne vais pas insister (\footnote[183]{J'ai vérifié que si $f$ est fibrant, alors $\cE$ est \textit{karoubienne}. Il y a donc de l'espoir \dots}).

\msp{page 439}Dans le cas où $f$ est fibrant, admettons (comme hypothèse supplémentaire au besoin) que $\cE$ satisfasse $L$ et que $f$ soit accessible, et prenons $\pi$ assez grand pour que $\cE$, $f$ satisfassent $\mathrm{L}_{\pi}$. Nous pouvons alors terminer la démonstration de 4.15.6. Il faut prouver
\[
\cE(\pi) \ \subset \ \cE_{\pi} \ ,
\]
donc que si $x \in \Ob \cE_{s\,\pi}$, et $(y_i)_{i \in I}$ est un système inductif dans $\cE$ grand devant $\pi$, de limite $y$, on a
\[
\varinjlim \Hom_{\cE}(x, y_i) \xrightarrow{\ \sim\ } \Hom_{\cE}(x, y) \ .
\]
Donc il faut prouver que toute flèche $x \xrightarrow{u} y$ se factorise, de façon essentiellement unique $u_i$, par un $y_i$. Mais la flèche $s \xrightarrow{\ \lambda \, = \, f(u)\ } t$ dans $\cC$ se factorise, de façon essentiellement unique, par des $\lambda_i \, : \, s \longrightarrow t_i$. Donc on est ramené à trouver un \textit{$\lambda_i$-morphisme} $x \longrightarrow y_i$ pour $i$ grand, qui factorise. Dans le cas fibrant, cela signifie qu'on cherche $x \longrightarrow \lambda^*_i(y_i)$ qui factorise $x \longrightarrow \lambda^*(y)$, or les $\lambda^*_i(y_i)$ forment un système inductif dans $\cE_s$, et sa limite \textit{dans} $\cE_s$ est $\lambda^*(y)$ (\textit{à vérifier}). Comme $x$ est dans $\cE_{s\,\pi}$, l'hom\texttt{[omorphisme]} donné $x \longrightarrow \lambda^*(y)$ se factorise de

\msp{page 440}façon essentiellement unique par un $\lambda^*_i(y_i)$.

Cette fois par contre, c'est le cas $f$ cofibrant qui pose problème. En fait, je n'arrive pas à m'en tirer.

Ainsi, on trouve bien un foncteur
\[
\Ind_{\pi}(\cE(\pi)) \ \longrightarrow \ \cE \ ,
\]

% ---- from der-ch18_p171-175.tex
\setcounter{footnote}{183}

\printedpage{169}
% head: A. Grothendieck : Dérivateurs, Ch. XVIII, Édition des Universités de Montpellier II et Paris VII
qui est essentiellement surjectif, mais faute de savoir que $\mathcal{E}(\pi) \subset \mathcal{E}_\pi$ dans le cas cofibrant, il n'est pas clair que ce foncteur soit une équivalence de catégories.

\textbf{4.15.7.} Il me faut reporter la démonstration du lemme 4.15.3 (p. 428), en admettant 4.8.13. (Non encore démontré -- je vais revenir plus bas sur la démonstration de 4.8.13.) On va utiliser le

\textbf{Lemme 4.15.7.1.} \textit{Soient $\cC$ une catégorie, $\ux = (x_i)_{i\in I}$ un ind-objet de $\cC$. Pour que $\ux$ soit essentiellement constant, i.e. isomorphe, dans} $\Ind(\cC)$\textit{, au ind-objet constant défini par un objet}

\msp{page 441}
\textit{de $\cC$, il faut et il suffit qu'il existe un foncteur cofinal $\varphi : J \longrightarrow I$, avec $J$ filtrant, de telle façon que le système réindexé $(x_j)_{j\in J}$, avec $x_j = x_{\varphi(j)}$, ait la propriété suivante : pour tout $j \in J$, on peut trouver un projecteur $p_j$ dans $x_j$, de telle façon que :}
\begin{itemize}
\item[\textit{1°)}] \textit{Pour toute flèche $\lambda : j \longrightarrow j'$ dans $J$, on ait commutativité dans}
\[
\begin{tikzcd}
x_j \arrow[r, "x_\lambda"] \arrow[d, "p_j"] & x_{j'} \arrow[d, "p_{j'}"] \\
x_j \arrow[r, "x_\lambda"] & x_{j'}
\end{tikzcd}
\]
\textit{(où les $x_\lambda$ sont les flèches de transition), i.e. $p_{j'}x_\lambda = x_\lambda p_j$.}
\item[\textit{2°)}] \textit{Pour tout $j \in \Ob J$, il existe une flèche $\lambda : j \longrightarrow j'$ dans $J$ telle que $p_{j'}x_\lambda = x_\lambda p_j = x_\lambda$.}
\item[\textit{3°)}] \textit{Pour tout $j \in \Ob J$, le projecteur $p_j$ dans $x_j$ admet un image $x_j^\natural$.} % [sic] un image -> une image, gender agreement
\item[\textit{4°)}] \textit{Pour toute flèche $\lambda : j \longrightarrow j'$ dans $J$, la flèche $x_\lambda^\natural : x_j^\natural \longrightarrow x_{j'}^\natural$ induite par $x_\lambda$ (grâce à 1°) est un isomorphisme.}
\end{itemize}
\textit{De plus, si on peut supposer $J$ de la forme $i_0\backslash I$, le foncteur cofinal $J \longrightarrow I$ étant le foncteur canonique $i_0\backslash I \longrightarrow I$ ($i_0$ étant un objet fixé convenable de $I$)} \texttt{[phrase incomplète]}

\msp{page 442}
Notons tout de suite deux corollaires.

\textbf{Corollaire 4.15.7.2.} \textit{Supposons que $\cC$ soit sans projecteurs. Alors $\ux = (x_i)_{i\in I}$ est un ind-objet essentiellement constant si et seulement s'il existe un foncteur cofinal, $J \longrightarrow I$, avec $J$ filtrant (\footnote{et on peut supposer $J = i_0\backslash I \xrightarrow{\text{can.}} I$}), de telle façon que le système inductif réindexé $(x_j)_{j\in J}$ soit isomorphe} (dans la catégorie $\underline{\mathrm{Hom}}(J,\cC)$) \textit{à un système inductif} constant \textit{(i.e. défini par un foncteur constant $J \longrightarrow \cC$ : tous les $x_j$ égaux à un même $x$, et les morphismes de transition des identités).}

La suffisance est claire, pour la nécessité on applique 4.15.7.1, qui nous dit qu'on peut réindexer de façon à satisfaire aux conditions 1°) à 4°) (pour des projecteurs $p_j$ convenables dans les $x_j$). Mais l'hypothèse sur $\cC$ implique que $p_j = \id_{x_j}\ \forall\, j \in \Ob J$, et les conditions 1°, 2°, 3° sont satisfaites tautologiquement pour un tel choix, la condition 4° signifie que les morphismes de transition sont des isomorphismes. Mais il est immédiat que pour un système inductif $(x_j)_{j\in J}$ donné, avec $J$ 1-\textit{connexe} (p.ex. $J$ filtrant), la condition que les $x_j$ sont des isomorphismes % [sic] x_j -> x_\lambda, the transition morphisms are meant
est nécessaire et suffisante pour que ce système soit isomorphe à un système inductif constant, d'où la conclusion.
% foot: Version 17 décembre 2025

\printedpage{170}
% head: A. Grothendieck : Dérivateurs, Ch. XVIII, Édition des Universités de Montpellier II et Paris VII
\msp{page 443}
\textbf{Corollaire 4.15.7.3.} \textit{Supposons que $I$ soit grande devant $\pi_0$, où $\pi_0$ est un cardinal infini tel que $\pi_0 \geqslant \card \cC$. Alors $(x_i)_{i\in I}$ est essentiellement constant si et seulement si on peut trouver un foncteur cofinal $J \longrightarrow I$ avec $J$ grande devant $\pi_0$ (et ordonnée si on veut) (\footnote{La démonstration (légèrement réarrangée) montre qu'on peut prendre pour $J$ une sous-catégorie (pas pleine) d'une catégorie $i_1\backslash I$, pour un objet convenable $i_1$ de $I$. On a donc $\card J \leqslant \pi$ si $\card I \leqslant \pi$, $\pi$ étant un cardinal infini donné.}), de telle façon que le système réindexé $(x_j)_{j\in J}$ satisfasse les conditions suivantes :}
\begin{itemize}
\item[\textit{1°)}] \textit{L'application $j \longmapsto x_j$ de $\Ob J$ dans $\Ob\cC$ est constante, i.e. il existe un $x_0 \in \Ob\cC$ tel que $x_j = x_0\ \forall\, j \in \Ob J$.}
\item[\textit{2°)}] \textit{Il existe un projecteur $p$ dans $x_0$ tel que pour $j < j'$ dans $J$, on ait $\tau_{j',j} = p$.}
\item[\textit{3°)}] \textit{Le projecteur $p$ admet une image dans $\cC$.}
\end{itemize}
\textsc{Démonstration}. Ces conditions impliquent évidemment celles de 4.15.7.1, en prenant $p_j = p\ \forall\, j \in \Ob J$, du moins si $J$ n'admet de plus grand élément (ce qui fait une difficulté pour 2°). Mais si $J$ admet un objet final, il en est de même de $I$, et il est clair alors que tout système inductif indexé par $I$ est essentiellement constant. Donc, dans ces cas, les conditions 1°) 2°) 3°) ci-dessus impliquent que $\ux$ est essentiellement constant. Inverse\msp{page 444}ment, supposons que $\ux$ soit essentiellement constant, et prouvons que l'on peut réindexer comme indiqué. Le cas où $I$ admet un objet final est trivial, on prendra $J = e$, $p = \id$. Supposons donc que $I$ n'ait pas d'objet final. Comme $I$ grande devant $\pi_0$ et $\pi_0 > \card\Ob\cC$, on voit qu'il existe un $x_0 \in \Ob\cC$ tel que la sous-catégorie pleine $I' = I_{x_0}$ de $I$, formée des $i$ tels que $x_i = x_0$, soit cofinale dans $I$. Comme $(x_i)_{i\in I}$ est essentiellement constant, il en est de même de $(x_{i'})_{i'\in I'}$, et on peut lui appliquer le critère 4.15.7.1, avec \marginpar{$\Vert$}% double vertical bar printed in the left margin beside the two following lines
$J \longrightarrow I'$ cofinal. La démonstration montrera de plus que si $I'$ est grand devant $\pi_0$ (cardinal infini sans autres conditions), on peut supposer $J$ grande devant $\pi_0$ (\footnote{en tenir compte dans la démonstration de 4.15.7.1.}). De plus, il est clair qu'on peut supposer $J$ ordonnée. Comme $\pi_0 \geqslant \card\Hom(x_0,x_0)$, on voit encore qu'il existe un projecteur $p$ dans $x_0$ tel que la partie $J_1 = J'$ de $J$, formée des $j$ tels que $p_j = p$, soit cofinale. Étant cofinale, elle est également grande devant $\pi_0$. Donc on peut supposer que tous les $p_j$ sont égaux à un même $p$. Les conditions 1°-4° deviennent :
\begin{itemize}
\item[1°)] $p\tau_{j'j} = \tau_{j'j}p$ pour $j \leqslant j'$.
\item[2°)] Pour tout $j$, existe $j' \geqslant j$ avec
\[
p\tau_{j'j} \;=\; \tau_{j'j}p \;=\; \tau_{jj'}\ . % [sic] \tau_{jj'} -> \tau_{j'j}, index order used everywhere else
\]
\item[3°)] L'image $x_0^\natural$ du projecteur $p$ dans $x_0$ existe.
\item[4°)] L'endomorphisme $\tau^\natural_{j'j}$ de $x_0^\natural$ est un isomorphisme, pour $j \leqslant j'$.
\end{itemize}
\msp{page 445}
Or on a le

\textbf{Lemme 4.15.7.4.} \textit{Soient $\cC$ une catégorie, $x_0$ un objet de $\cC$, $p$ un projecteur dans $x_0$, $\tau$ un endomorphisme de $x_0$ satisfaisant les relations}
\[
\tau p \;=\; p\tau \;=\; \tau\ , \tag*{$(*)$}
\]
% foot: Version 17 décembre 2025
\printedpage{171}
% head: A. Grothendieck : Dérivateurs, Ch. XVIII, Édition des Universités de Montpellier II et Paris VII
\textit{et tel que l'automorphisme de $x_0^\natural \defeq$} $\Img p$ \textit{(défini p.ex. comme objet de $\cC^\wedge$, $x_0^\natural =$} $\Ker(\id_{x_0},p)$ % [sic] "(défini ... Ker(id_{x_0},p)" -> a closing parenthesis is missing, the parenthesis opened before "défini" is never closed
\textit{soit l'identité. Alors on a $\tau = p$.}

\textsc{Démonstration de} 4.15.7.4. On est ramené aussitôt, par le procédé standard, au cas où $\cC = \mathrm{Ens}$, où la vérification est immédiate : les automorphismes $\tau$ de l'ensemble $E$ $(= x)$ qui satisfait $(*)$ correspondent biunivoquement aux endomorphismes $u$ de $E^\natural = pE$, via
\[
\begin{cases}
u \;\longmapsto\; \sigma u \rho \\
\tau \;\longmapsto\; \tau|E^\natural
\end{cases}
\qquad
\text{(où }
\begin{tikzcd}[column sep=small]
E^\natural \arrow[r, hook, shift left, "\sigma"] & E \arrow[l, shift left, "\rho"]
\end{tikzcd}
\text{ sont les flèches canoniques)}\ ,
\]
donc $\tau|E^\natural = \id \Longrightarrow u = \sigma\rho = p$.

Choisissons un $j_0 \in J$, et considérons les homomorphismes
\[
\tau^\natural_{j,j_0} \;:\; x_0^\natural \longrightarrow x_0^\natural
\]
induits par les $\tau_{j,j_0}$ avec $j \geqslant j_0$. Ici encore, comme l'ensemble des isomorphismes possibles est de

\msp{page 446}
cardinal $\leqslant \pi_1$, et $J/j_0$ est grande devant $\pi_0$, on voit qu'il existe $J' \subset j_0\backslash J$, cofinale dans $j_0\backslash J$, donc dans $J$, telle que pour $j \in J'$, les $\tau^\natural_{j,j_0}$ soient égaux à un même $\tau : x_0^\natural \xrightarrow{\sim} x_0^\natural$. Donc si $j', j'' \in J'$, avec $j' \leqslant j''$, on a
\[
\underbrace{\tau_{j''j_0}}_{\tau} \;=\; \tau_{j''j'}\,\underbrace{\tau_{j'j_0}}_{\tau}\ ,
\]
d'où
\[
\tau_{j''j'} \;=\; \id_{x_0^\natural}\ ,
\]
puisque la flèche $\tau$ est un isomorphisme. Donc le système inductif réindexé par $J'$, la condition 4° est remplacée par la condition plus stricte \fbox{4° bis) $\tau^\natural_{j',j} = \id_{x_0^\natural}$.} Mais la condition 2° devient alors, grâce à 4.15.7.4,
\[
\boxed{\text{2° bis) Pour tout } j\text{, existe } j' \geqslant j \text{ avec } \tau_{j',j} = p}\ .
\]
Donc, on peut supposer 1°, 2° bis, 3°, 4° bis (on désigne à nouveau par $J$, au lieu de $J'$, l'ensemble ordonné grand devant $\pi_0$ indexant).

Soit $K$ l'ensemble ordonné défini par

\msp{page 447}
\[
\begin{cases}
\Ob K \;=\; \Ob J \\
\text{si } j,j' \in \Ob J\text{, on a } j \overset{K}{\leqslant} j' \text{ si et seulement si } j \overset{\mathrm{J}}{=} j'\text{, ou } j \overset{J}{<} j' \text{ et } \tau_{j',j} = p\,.
\end{cases}
\]
Il est immédiat que c'est une \textit{relation d'ordre} sur $J$ (car $p^2 = p$), plus fin que la relation d'ordre initiale. Le système inductif induit sur $K$ satisfait aux conditions 1°, 2°, 3° de 4.15.7.3. D'autre part, $K \longrightarrow J$ est cofinal, donc aussi le composé $K \longrightarrow J \longrightarrow I$. Il reste à prouver que $K$ est grande devant $\pi_0$, donc que toute partie $K' \subset K$ telle que $\card K' \leqslant \pi_0$ admet un majorant dans $K$. Elle admet un majorant \textit{dans} $J$, soit $j$. Par hypothèse 2 bis, il existe $j' \geqslant j$ tel que $\tau_{j'j} = p$, et on peut supposer $j' > j$. (En effet, si $\tau_{j'j} = p$, on voit que
% foot: Version 17 décembre 2025
\printedpage{172}
% head: A. Grothendieck : Dérivateurs, Ch. XVIII, Édition des Universités de Montpellier II et Paris VII
\msp{page 448}
la forme initiale 2°) de la condition 2° bis que l'on a $\tau_{j''j} = p$ pour tout $j'' \geqslant j'$, donc si $j$ n'est par objet final de $J$, % [sic] par -> pas, misprint
on peut supposer $j' > j$. Mais $J$ n'a pas d'objet final, car sinon $I$ en aurait un, cas que l'on avait éliminé. Donc il existe $j'$ tel que $j' \overset{K}{>} j$.) Or on a ceci (qu'on applique au cas $i \in K'$) :
\[
i \leqslant j \ \text{ et } \ j \overset{K}{<} j' \qquad \Longrightarrow \qquad i \overset{K}{<} j'\ .
\]
\[
\underbrace{x_i}_{= x_0} \xrightarrow{\tau_{j,i}} x_j \xrightarrow{\tau_{j',j}=p} x_{j'}\ ,
\]
\[
\tau_{j'i} \;=\; \tau_{j'j}\tau_{ji} = \underbrace{p\tau_{ji}}_{\overset{\text{1°}}{=}\,\tau_{ji}p} \;\overset{?}{=}\; p\ ,
\]
on a en effet
\[
\left.
\begin{aligned}
\tau_{j',i} &\;=\; p\tau_{j',i} \\
\tau_{j',i} &\;-\; \tau_{j',i}p % [sic] - -> =, the second relation of (*) is meant
\end{aligned}
\right\}
\quad \overset{\text{par 4.15.7.4}}{\Longrightarrow} \quad % Doubt resolved against PDF leaf 14: no underline; the two strokes belong to the implication arrow.
\tau_{j',i} \;=\; p\ .
\]
Cela achève la démonstration du corollaire 4.15.7.3.

\textbf{Corollaire 4.15.7.5} (de 4.15.7.3). \textit{Soit $(x_i)_{i\in I}$ un système inductif essentiellement constant grand devant $\pi_0 \geqslant$} $\card\cC$\textit{, dans $\cC$} sans projecteurs. \textit{Alors il existe $J \longrightarrow I$ cofinal, avec $J$ ordonné grand devant $\pi_0$, et tel que le système inductif induit sur $J$ soit} constant.

\textsc{Démontrons le lemme} 4.15.7.1. Il faut expliciter l'existence de deux flèches dans $\Ind(\cC)$,
\[
z \underset{\beta}{\overset{\alpha}{\rightleftarrows}} \text{"}\varinjlim_{i\in I}\text{"} x_i \;=\; \ux\ ,
\]
avec $\beta\alpha = \id_z$, $\alpha\beta = \id_{\ux}$. La donnée

\msp{page 449}
de $x$ équivaut à celle d'un élément de $\varinjlim\limits_{i\in I} \Hom(z,x_i)$, et est donc donné par un % [sic] x -> \alpha, parallel with "la donnée de \beta" below
\[
\alpha_{i_0} \;:\; z \longrightarrow x_{i_0}\ ,
\]
la donnée de $\beta$ équivaut à celle d'un élément de $\varprojlim\limits_{i\in I} \Hom(x_i,z)$, i.e. de flèches
\[
\beta_i \;:\; x_i \longrightarrow z\ ,
\]
avec la condition de compatibilité habituelle. Quitte à remplacer $I$ par $I'$, avec $I' \longrightarrow I$ cofinal, on peut supposer $I$ ordonnée (et grande devant $\pi_0$, si $I$ initiale l'était), et la condition de compatibilité sur les $\beta_i$ s'écrit
\[
\beta_i\tau_{ji} \;=\; \beta_j \qquad \text{si } i \leqslant j \text{ dans } I,
\]
où $\tau_{ji} : x_i \longrightarrow x_j$ désigne la flèche de transition. Il faut expliciter, en termes de ces données, les conditions
\[
\alpha\beta \;=\; \id_z\ , \qquad \beta\alpha \;=\; \id_{\ux}\ .
\]
% foot: Version 17 décembre 2025

\printedpage{173}
% head: A. Grothendieck : Dérivateurs, Ch. XVIII, Édition des Universités de Montpellier II et Paris VII
\[
\begin{tikzcd}
z \arrow[r, "\alpha_{i_0}"] \arrow[dr] & x_{i_0} \arrow[r, "\tau_{i,i_0}"] \arrow[d, "\beta_{i_0}"] & x_i \arrow[dl, "\beta_i"] \\
 & z &
\end{tikzcd}
\]
La relation $\alpha\beta = \id_z$ équivaut à
\[
\beta_{i_0}\alpha_{i_0} \;=\; \id_z\ , \tag*{$(*_{i_0})$}
\]
elle implique d'ailleurs, en posant pour $i \geqslant i_0$,
\[
\alpha_i \;\defeq\; \tau_{i,i_0}\alpha_{i_0} \;:\; z \longrightarrow x_i\ ,
\]
que l'on a pour $i \geqslant i_0$
\[
\beta_i\alpha_i \;=\; \id_z \qquad \forall\, i \geqslant i_0\ . \tag*{$(*_i)$}
\]

\msp{page 450}
Posons alors
\[
p_i \;=\; \alpha_i\beta_i \;:\; x_i \longrightarrow x_i\ ,
\]
il est clair qu'on a
\[
\begin{cases}
p_i^2 \;=\; p_i \qquad \text{($p_i$ un projecteur)} \\
\alpha_i \;:\; z \longrightarrow x_i \qquad \text{induit un isomorphisme} \\
\qquad z \xrightarrow{\sim} x_i^\natural \;\defeq\; \Img p_i \;=\; \Ker(\id_{x_i},p_i)\ .
\end{cases}
\]
L'ensemble ordonné cofinal cherché $J$ est $i_0\backslash I$, il est filtrant, et est grand devant $\pi_0$ si $I$ l'est. À vrai dire, la réduction au cas $I$ ordonné est inutile, on peut prendre directement $J = i_0\backslash I$, il est immédiat que $J$ est filtrante ($I$ l'étant), et $i_0\backslash I \longrightarrow I$ cofinal, enfin que si $I$ est grande devant un cardinal infini $\pi_0$, il en est de même de $i_0\backslash I$.

Exprimons que $\beta\alpha = \id_{\ux}$ ; or on voit que $\beta\alpha$ est défini par l'endomorphisme $(p_j)_{j\in J}$ du système inductif $(x_j)_{j\in J}$ -- le fait que ce soit bien un endomorphisme de ce système, i.e. qu'on ait la relation 1°) de 4.15.7.1,
\[
p_{j'}x_\lambda \;=\; x_\lambda p_\lambda \qquad \text{si } \lambda : j \longrightarrow j' \text{ dans } J, % [sic] p_\lambda -> p_j, cf. relation 1°) of 4.15.7.1
\]
est immédiat, ça s'écrit
\[
\alpha_{j'}\underbrace{\beta_{j'}x_\lambda}_{\beta_j} \;=\; \underbrace{x_\lambda\alpha_j}_{\alpha_{j'}}\beta_j\ , \qquad \text{or les deux membres sont égaux à } \alpha_{j'}\beta_j.
\]
% foot: Version 17 décembre 2025

% ---- from der-ch18_p176-180.tex
\printedpage{174}
% head: A. Grothendieck : Dérivateurs, Ch. XVIII, Édition des Universités de Montpellier II et Paris VII

\msp{page 451}
Dire que c'est l'identité, est exprimé par la condition 2°) sur $(x_j,p_j)_{j\in J}$. La condition 3° est clair par construction, $p_j$ admettant une image isomorphe à $z$ par $\alpha_j$. D'autre part, on voit \texttt{[tout]} de suite que pour une flèche $\lambda : j \longrightarrow j'$ de $J$, on a commutativité dans % [sic] est clair -> est claire, agreement with "condition"
\[
\begin{tikzcd}
x_j \arrow[r, "x_\lambda"] & x_{j'} \\
x_j^\natural \arrow[u, hook] \arrow[r, "x_\lambda"] & x_{j'}^\natural \arrow[u, hook] \\
z \arrow[u, "\wr", "\overline{\alpha_j}"'] \arrow[r, "\mathrm{id}_z", "\sim"'] & z \ , \arrow[u, "\wr", "\overline{\alpha_{j'}}"']
\end{tikzcd}
\]
où $\overline{\alpha_j}$, $\overline{\alpha_{j'}}$ sont les isomorphismes définis par $\alpha_j$, $\alpha_{j'}$ respectivement. Cela prouve que $x_\lambda^\natural$ est un isomorphisme.

Cela prouve que si $(x_i)_{i\in I}$ est essentiellement constant, il existe $J\longrightarrow I$ comme indiqué. L'inverse est également vrai, car on peut supposer $J = I$, puisque $(x_i)_{i\in I}$ et $(x_j)_{j\in J}$ définissent des ind-objets isomorphes. Alors le système inductif des $x_j^\natural$ a des morphismes de transition des isomorphismes, et comme $J$ est filtrante donc 1-connexe, on en déduit facilement qu'il est isomorphe
\msp{page 452}
à un système inductif constant, de valeur $z$. En choisissant un tel isomorphisme
\[
z \xrightarrow[\sim]{\overline{\alpha_i}} x_i^\natural \ ,
\]
on trouve $z \xrightarrow{\alpha} \underline{x}$, et le système des $p_i$ définit des
\[
\begin{tikzcd}
x_i \arrow[r, "\overline{p_i}"] \arrow[rr, bend right=12, "\beta_i \overset{\text{déf}}{=} \overline{\alpha_j}^{-1}\overline{p_i}"'] & x_i^\natural \arrow[r, "\overline{\alpha_j}^{-1}"] & z \ ,
\end{tikzcd}
\] % [sic] \overline{\alpha_j} (index j) -> \overline{\alpha_i}, the arrow goes from x_i^\natural, cf. z -> x_i^\natural above
et on voit facilement que l'on aura
\[
\beta\alpha \ = \ id_z \ , \qquad \alpha\beta \ = \ \mathrm{id}_{\underline{x}} \ ,
\]

\hfill q.e.d.

\textsc{Je peux maintenant prouver} 4.15.3. C'est en effet essentiellement la conjonction des 4.15.7.5 (p. 448), et de 4.8.13, qui nous dit essentiellement que si $\mathcal{C}$ est \textit{karoubienne} (condition vérifiée si $\mathcal{C}$ est sans projecteurs), alors tout système inductif dans $\mathcal{C}$, grand devant $\pi_0$, où $\pi_0 \geqslant \card\,\text{réd}\,\mathcal{C}$, est essentiellement constant,

\hfill q.e.d.

\printedpage{175}
% head: A. Grothendieck : Dérivateurs, Ch. XVIII, Édition des Universités de Montpellier II et Paris VII

\msp{page 453}
\textbf{4.15.8.} Je veux revenir ici sur la démonstration inachevée de 4.8.13 (p. 212-216). Je reprends la démonstration où je l'avais laissée (p. 215). Il faut prouver que
\[
F(\xi) \ = \ \varinjlim_{i\in I} M_i
\]
est de cardinal $\leqslant \pi_0$. J'ai annoncé qu'on peut se ramener au cas où $M$ est un monoïde de monomorphismes, i.e. que $uv = uv' \Longrightarrow v = v'$ (pour $u,v,v' \in M$). Je vais expliciter cet argument.

Pour tout $i\in I$, soit
\[
N_i \ = \ \left\{ u \in M \ \middle| \ \begin{array}{l}
\text{l'ensemble } I_{i,u} \text{ des } j\in I \text{ tels}\\
j \geqslant i,\ \alpha_{ji} = u \text{ est cofinal}\\
\text{dans } I\text{, i.e. pour tout } j_0\in I,\\
\text{existe } j \in I_{i,u}\text{, avec } j \geqslant j_0
\end{array} \right\} .
\tag*{(4.15.8.1)}
\]
Comme $I$, donc $i\backslash I$ est grande devant $\pi_0$, et $\card M \leqslant \pi_0$, on voit aussitôt que
\[
N_i \ \neq \ \emptyset \ .
\tag*{(4.15.8.2)}
\]
$\Big($ Si d'autre part $i' \geqslant i$, on trouve
\[
N_i \ = \ N_{i'}\alpha_{i'i} \ .
\tag*{(4.15.8.3)}
\]
L'inclusion $N_{i'}\alpha_{i'i} \subset N_i$ est $\pm$ tautologique, prouvons $N_i \subset N_{i'}\alpha_{i'i}$. Soit $u \in N_i$, l'ensemble $J = I_{i,u} \cap (i'\backslash I)$ des $j \geqslant i'$ tels que $\alpha_{ji} = u$ est cofinal dans $I$, donc grand devant $\pi_0$.
\[
\begin{tikzcd}
i \arrow[r] \arrow[rr, bend right=12, "u = \alpha_{ji}"'] & i' \arrow[r, "\alpha_{ji'}"] & j
\end{tikzcd}
\]
Donc il existe un $v \in M$ tel que l'ensemble $J'$ de $j\in J$ tels que $\alpha_{ji'} = v$
\msp{page 454}
soit cofinal dans $J$, donc dans $I$. Comme cet ensemble $J'$ est contenu dans $I_{i',v}$, cela implique que $I_{i',v}$ est cofinal dans $I$, i.e. $v \in N_{i'}$. On a d'autre part, pour $j \in J'$,
\[
\underbrace{\alpha_{ji'}}_{v}\alpha_{i'i} \ = \ \underbrace{\alpha_{ji}}_{u} \ ,
\]
d'où $u = v\alpha_{i'i}$ avec $v\in N_{i'}$, d'où $u \in N_i\alpha_{i'i}$, q.e.d. $\Big)$ % [sic] N_i\alpha_{i'i} -> N_{i'}\alpha_{i'i}, the inclusion to be proved is N_i \subset N_{i'}\alpha_{i'i}

Pour $i \in I$, soit
\[
\alpha_i \ : \ \underbrace{\mathrm{Hom}(\xi, \underbrace{x_i}_{=\,\xi})}_{=\,M} \ \longrightarrow \ \underbrace{\mathrm{Hom}(\xi,F) \ = \ F(\xi)}_{\simeq\ \varinjlim\limits_{i} \mathrm{Hom}(\xi,x_i)}
\]
\printedpage{176}
% head: A. Grothendieck : Dérivateurs, Ch. XVIII, Édition des Universités de Montpellier II et Paris VII
le morphisme canonique. On trouve ceci :
\begin{center}
\makebox[0.20\linewidth][l]{(4.15.8.4)}%
\begin{minipage}[c]{0.78\linewidth}
Si $u \in N_i$, alors la relation d'équivalence $R_i$ dans $M$ définie par $\alpha_i$ ($v \overset{R_i}{\sim} v' \Longleftrightarrow \alpha_i(v) = \alpha_i(v')$) est égale à celle définie par l'application $v \longmapsto uv$ de $M$ dans $M$, i.e.
\[
\alpha_i(v) \ = \ \alpha_i(v') \qquad \Longleftrightarrow \qquad uv \ = \ uv' \ .
\]
\end{minipage}
\end{center}

C'est immédiate, puisque $I_{i,u}$ étant cofinal dans $I$, on a $\alpha_i(v) = \alpha_i(v')$ si et seulement si $\exists\ j \in I_{i,u}$ avec $\alpha_{ji}v = \alpha_{ji}v'$, or $\alpha_{ji} = u$. % [sic] C'est immédiate -> C'est immédiat, agreement

Pour $u \in M$, notons $R_u$ la relation d'équivalence dans $M$ définie par $v \longmapsto uv$. L'ensemble des relations d'équivalence dans \texttt{[}$M$\texttt{]}
\msp{page 455}
obtenu ainsi est de cardinal $\leqslant \card M \leqslant \pi_0$, donc l'ensemble des relations d'équivalences $R_i$ ($i \in I$) est lui aussi de cardinal $\leqslant \pi_0$. Cela montre, $\pi_0$ étant grand devant $\pi_0$, qu'il existe une relation d'équivalence de la forme $R_u$ ($u \in M$), et un ensemble cofinal $I'$ dans $I$, tel que $\forall\ i \in I'$ on ait $R_i = R_u$. Cela montre que, quitte à remplacer $I$ par l'ensemble ordonné cofinal $I'$, qu'on peut supposer que l'on a
\[
R_i \ = \ R_u \qquad \forall\, i \in I \ .
\tag*{(4.15.8.5)}
\]
Soit alors
\[
E \ = \ M/R_u \ = \ M/R_i \qquad (\forall\ i \in I) \ ,
\]
on voit donc que les $\alpha_{ji} \in M$ (pour $i \leqslant j$) sont tels que l'application
\[
\tau_{\alpha_{ji}} \ : \ v \ \longmapsto \ \alpha_{ji}v \ : \ M \ \longrightarrow \ M
\]
est compatible avec la relation d'équivalence $R_u$, donc passe au quotient en une application
\[
\overline{\alpha_{ji}} \ : \ E \ \longmapsto \ E \ ,
\]
laquelle est d'ailleurs \textit{injective.}

Soit $M'$ le sous-monoïde de $M$ formé des $u \in M$ qui ont la propriété de stabilité précédente, ainsi que la propriété d'injectivité qu'on vient de dire, et $\overline{M'}$ le monoïde image de $M'$ dans le monoïde
\[
\mathrm{Mon}(E) \ \subset \ \mathrm{End}_{\mathrm{Ens}}(E)
\]
\msp{page 456}
des applications injectives de $E$ dans lui-même. On voit donc que le système inductif donné, défini par les $\alpha_{ji}$, prend ses valeurs dans le sous-monoïde $M'$ de $M$ (i.e. de la sous-catégorie correspondante de la catégorie $B_M$), et que la limite inductive $F(\xi)$ qui nous intéresse (pour prouver $\card F(\xi) \leqslant \pi_0$) est isomorphe à la limite inductive du système inductif d'un système d'ensembles tous égaux à $E$, les morphismes de transition étant les $\overline{\alpha_{ji}} \in \overline{M'} \subset \mathrm{Mon}(E)$. De plus, on aura $\card \overline{M'} \leqslant \card M \leqslant \pi_0$. Le résultat $\card F(\xi) \leqslant \pi_0$ voulu résultera alors du

\textbf{Lemme-clef 4.15.8.1.} \textit{Soit $M$ un monoïde de monomorphismes (i.e. tel que $uv = uv'$ $\Longrightarrow v = v'$, pour $u,v,v' \in M$), et soit $(\alpha_{ji})_{i\leqslant j}$ un système inductif dans $B_M$, indexé par un ensemble ordonné $I$ grand devant $\pi_0$, où $\pi_0$ est un cardinal}
\printedpage{177}
% head: A. Grothendieck : Dérivateurs, Ch. XVIII, Édition des Universités de Montpellier II et Paris VII
\msp{page 457}
\textit{infini tel que} $\card M \leqslant \pi_0$. \textit{Alors il existe un $i_0 \in I$ tel que pour tout $i \geqslant i_0$, $\alpha_{i_0,i}$ soit un isomorphisme.}

Si on applique ce lemme au cas de $\overline{M'}$ et du $\overline{\alpha_{ji}}$, on trouve donc qu'il existe un $i_0 \in I$ tel que les $\overline{\alpha_{ji_0}}$ soient des isomorphismes. Cela implique que $F(\xi) \simeq \varinjlim\limits_{i} \underbrace{E_i}_{E}$ est isomorphe à $E_{i_0}$, donc est de cardinal égal à $\card E_{i_0} = \card E \leqslant \pi_0$, ce qui achèvera la démonstration.

D'ailleurs, l'énoncé 4.8.13 implique 4.15.8.1, car dans le système inductif des $(M_i = M, \alpha_{ji})$, de $\varinjlim$ égale à $F(\xi)$, s'il n'existait \textit{pas} de $i_0$ comme stipulé dans 4.15.8.1, alors on en conclurait aisément que l'on aurait $\card F(\xi) > \pi_0$, contrairement à 4.8.13. Nous allons dégager ceci dans un lemme :

\msp{page 458}
\textbf{Lemme 4.15.8.2.} \textit{Soit $(E_i)_{i\in I}$ un système inductif d'ensembles, indexé par un ensemble ordonné $I$ grand devant $\pi_0$. On suppose qu'il existe $i_0 \in I$ tel que}
\begin{enumerate}[nosep]
\item[\textit{a)}] \textit{$i \geqslant i_0$ implique} $\card E_i = \pi_0$, \textit{et}
\item[\textit{b)}] \textit{$j \geqslant i \geqslant i_0$ implique que $u_{ji} : E_i \longrightarrow E_j$ est un monomorphisme. Supposons de plus}
\item[\textit{c)}] \textit{$\forall\ i \geqslant i_0$ $\exists\ j > i$ tel que $u_{ji}$ ne soit} pas \textit{un isomorphisme (i.e. ne soit pas surjectif). Soit $E_\infty = \varinjlim\limits_{i\in I} E_i$. Alors} $\card E_\infty > \pi_0$.
\end{enumerate}
Ce lemme implique donc l'\textit{équivalence} de 4.8.13 et du “lemme-clef” 4.15.8.1. Pour prouver 4.15.8.2, nous allons prouver quelque chose d'un peu plus précis :

\textbf{Corollaire 4.15.8.3.} \textit{Sous les conditions de 4.15.8.2, soit $\alpha_0$ le premier ordinal tel que} $\card \alpha_0 > \pi_0$, \textit{et soit $I_{\alpha_0}$ l'ensemble des ordinaux $\beta < \alpha$ (ensemble ordonné qui est grand devant $\pi_0$). Il existe alors une application strictement croissante}
\[
I_{\alpha_0} \ \longrightarrow \ I \ , \qquad \alpha \ \longmapsto \ i_\alpha \ , \qquad \textit{appliquant } 0 \textit{ dans } i_0 \ ,
\]
\msp{page 459}
\textit{tel que pour tout $\alpha \in I_{\alpha_0}$, on ait}
\[
\underbrace{u_{i_{\alpha+1},i_\alpha}}_{\substack{\textit{c'est un } \text{monomorphisme,}\\ \textit{d'après b, vu que } i_{\alpha+1} > i_\alpha \geqslant i_0}} \qquad\qquad \text{\textit{n'est pas un isomorphisme} (\footnotemark[187]).}
\]
\footnotetext[187]{\textit{mieux} : $\forall\ \alpha \in I_{\alpha_0}$ $\varinjlim\limits_{\beta\in I_\alpha} E_{i_\beta} \longrightarrow E_{i_\alpha}$ n'est pas surjectif.}
\textit{On a alors}
\[
E'_\infty \ \overset{\text{déf}}{=} \ \underbrace{\varinjlim_{\alpha\in I_{\alpha_0}} E_{i_\alpha}}_{} \ \hookrightarrow \ \varinjlim_{i\in I} E_i \ \overset{\text{déf}}{=} E_\infty
\]
\textit{et}
\[
\card E'_\alpha \ \geqslant \ \underbrace{\card I_{\alpha_0}}_{=\ \textit{successeur de } \pi_0} \ > \ \pi_0 \ ,
\] % [sic] E'_\alpha -> E'_\infty, the following line has \card E'_\infty
\textit{d'où} $\card E_\infty \geqslant \card E'_\infty > \pi_0$.

\textsc{Démonstration}. La construction des $i_\alpha$ par récurrence transfinie est immédiate, en utilisant c) et le fait que $I$ est grande devant $\pi_0$, et que pour tout $\alpha \in I_{\alpha_0}$, l'ensemble $I_\alpha$
\printedpage{178}
% head: A. Grothendieck : Dérivateurs, Ch. XVIII, Édition des Universités de Montpellier II et Paris VII
des $\beta < \alpha$ est de cardinal $\leqslant \pi_0$. Le fait que $E'_\infty \longrightarrow E_\infty$ soit injectif résulte de b). Le fait que
\[
\card E'_\infty \ \geqslant \ \card I_{\alpha_0}
\]
résulte du fait que $E'_\infty$ est réunion croissante de sous-ensembles identifiés à $E_{i_\alpha}$, et pour tout
\msp{page 460}
$\alpha$ on peut trouver un
\[
\xi_\alpha \ \in \ E_{i_\alpha}\backslash \bigcup_{\beta<\alpha} E_{i_\beta} \ .
\]
L'application
\[
\begin{aligned}
\alpha \ &\longmapsto \ \xi_\alpha\\
I_{\alpha_0} \ &\longrightarrow \ E'_\infty
\end{aligned}
\]
est injective, donc $\card E'_\infty \geqslant \card I_{\alpha_0}$.

% vertical bar printed in the left margin beside the whole of Corollaire 4.15.8.4
\textbf{Corollaire 4.15.8.4.} \textit{Pour un $M$ et un $\pi_0$ donnés, pour que le lemme-clef 4.15.8.1 soit valable, il faut et il suffit qu'il le soit dans le cas où $I = I_{\alpha_0}$, $\alpha_0$ étant le premier ordinal tel que} $\card \alpha_0 > \alpha_0$. \textit{Il revient au même qu'on ne puisse trouver de système inductif $(\alpha_{ji})_{j\geqslant i}$ indexé par $I_{\alpha_0}$, à valeurs dans $M$, tel que pour $i < j$, on ait $\alpha_{ji}$ non inversible dans $M$.} % [sic] \card \alpha_0 > \alpha_0 -> \card \alpha_0 > \pi_0, cf. Corollaire 4.15.8.3

\msp{page 461}
Je vais résumer ce qui a été obtenu dans un énoncé récapitulatif :

\textbf{Théorème 4.15.9.} \textit{Soient $\mathcal{C}$ une petite catégorie, $\pi$ un cardinal tel que $\pi \geqslant$} $\card\mathcal{C}$. \textit{Les conditions suivantes sont équivalentes (et je conjecture qu'elles sont toujours satisfaites) :}
\begin{enumerate}[nosep]
\item[\textit{a)}] \textit{$\mathcal{M} \overset{\text{déf}}{=} \mathrm{Ind}_\pi\mathcal{C}$ est réduit à $\mathcal{M}_\pi$.}
\item[\textit{a$'$)}] \textit{(Si $\mathcal{C}$ est karoubienne.) Tout système inductif $\underline{x} = (x_i)_{i\in I}$ dans $\mathcal{C}$, grand devant $\pi$, est essentiellement constant.}
\item[\textit{a$''$)}] \textit{Pour tout système inductif $(x_i)_{i\in I}$ dans $\mathcal{C}$, grand devant $\pi$, il existe un foncteur cofinal $J \longrightarrow I$, avec $J$ ensemble ordonné grand devant $\pi$, et tel que le système réindexé $(x_j)_{j\in J}$ satisfasse les conditions suivantes :}
\begin{enumerate}[nosep]
\item[\textit{1)}] \textit{Tous les $x_j$ sont égaux à un même élément $x$ de $\mathcal{C}$.}
\item[\textit{2)}] \textit{Il existe un projecteur $p$ dans $x$, tel que l'on ait $\alpha_{j'j} = p$ pour $j < j'$ ($j,j' \in J$), où $\alpha_{j'j} : x_j \longrightarrow x_{j'}$ est la flèche de transition.}
\end{enumerate}
\item[\textit{b)}] \textit{Pour tout $x \in \mathrm{Ob}\,\mathcal{C}$, et pour tout monoïde quotient $M'$ de $M = \mathrm{End}_{\mathcal{C}}(x)$, tel que $M'$ soit un monoïde de monomorphismes (i.e. $uv = uv' \Longrightarrow v = v'$ si $u,v,v' \in M'$), et pour tout système inductif grand devant $\pi$ à valeurs dans $B_{M'}$ (définie par un système d'éléments $(\alpha_{ji})_{i\leqslant j,\ i,j\in I}$ d'un ensemble ordonné $I$ grand devant $\pi$, avec $\alpha_{kj}\alpha_{ji} = \alpha_{ki}$ pour $i \leqslant j \leqslant k$), il existe $i_0 \in I$ tel}
\msp{page 462}
\textit{que $\alpha_{ii_0} = 1$ pour $i \geqslant i_0$, i.e. $\alpha_{ji} = 1$ pour $j \geqslant i \geqslant i_0$.}
\item[\textit{b$'$)}] \textit{$M'$ étant un monoïde de monomorphismes comme dans b) (donc} $\card M' \leqslant \card\mathcal{C} \leqslant$ \textit{$\pi$), soit de plus $\alpha_0$ le plus petit ordinal tel que} $\card\alpha_0 > \pi$, \textit{et soit $I = I_{\alpha_0}$ l'ensemble ordonné des ordinaux $\alpha < \alpha_0$ (lequel ensemble est grand devant $\pi$). Il n'existe} pas \textit{de système inductif dans $B_{M'}$, tel que pour $i < j$ dans $I$, $\alpha_{ji}$ soit un élément non inversible dans $M'$ (donc $u \longmapsto \alpha_{ji}u : M' \longrightarrow M'$ est injectif mais non surjectif).}
\end{enumerate}

% ---- from der-ch18_p181-185.tex
\printedpage{179}
% head: A. Grothendieck : Dérivateurs, Ch. XVIII, Édition des Universités de Montpellier II et Paris VII
% foot: Version 17 décembre 2025

Signalons tout de suite le

\textbf{Corollaire 4.15.10.} \textit{Les conditions équivalentes précédentes sont satisfaites, si $\pi > \card\mathcal{C}$.}

En effet, soit $\alpha_1$ le plus petit ordinal tel que $\card\alpha_1 > \card\mathcal{C}$. Donc $\pi_1 \defeq \card\alpha_1$ est le successeur de $\pi_0 \defeq \card\mathcal{C}$, et on a $\pi \geqslant \pi_1$, donc (si $\alpha_0$ est comme explicité dans b$'$ ci-dessus) $\card\alpha_0 > \pi_1$, donc $\alpha_1 < \alpha_0$. Soit (avec les notations de b$'$), pour $\alpha \in I$ \texttt{[plutôt} $\alpha \in I_{\alpha_0}$\texttt{]}, i.e.\ $\alpha < \alpha_0$,
\[
M_\alpha \;=\; \varinjlim_{i<\alpha} \underbrace{M_i}_{=\,M}\ ,
\]
on voit donc, comme les $\alpha_{ji}$ sont injectifs et non surjectifs, que
\[
\card M_\alpha \;\geqslant\; \card\alpha\ ,
\]
\msp{463}
et que pour $\alpha < \beta$, on a
\[
M_\alpha \;\into\; M_\beta \qquad \text{injectif}\ .
\]
Appliquons ceci au cas où $\alpha = \alpha_1$, donc on trouve
\[
\card M_{\alpha_1} \;\geqslant\; \card\alpha_1 \;=\; \pi_1\ ,
\]
d'autre part, en prenant $\beta = \alpha_1 + 1$, d'où $M_\beta \simeq M$, on trouve une injection
\[
M_{\alpha_1} \;\into\; M\ ,
\]
d'où
\[
\pi_1 \;=\; \card M_{\alpha_1} \;\leqslant\; \card M \;=\; \pi_0\ ,
\]
une contradiction, ce qui établit b$'$.

\textbf{Corollaire 4.15.11.} \textit{Soit $\mathcal{C}$ une catégorie essentiellement petite, $\pi$ un cardinal tel que $\pi > \card\redop\mathcal{C}$. Soit $\mathcal{M} \defeq \Ind_\pi\mathcal{C}$, alors $\mathcal{M} = \mathcal{M}_\pi$, et c'est une enveloppe de Karoubi de $\mathcal{C}$.}

On peut supposer $\mathcal{C}$ réduite, donc $\card\mathcal{C} < \pi$, et on applique le corollaire précédent.

Ainsi, le résultat litigieux 4.8.13 (p.~212) est finalement prouvé sous une forme légèrement plus faible, avec $\pi > \card\redop\mathcal{C}$ au lieu de $\pi \geqslant \card\redop\mathcal{C}$. Du même coup, les applications de 4.8.13 qu'on a fait jusqu'à présent (notamment le lemme 4.15.3).

\msp{464}

\subsection*{4.16\quad Exemples et questions}

\textbf{4.16.0.} J'avais cru que tout foncteur “raisonnable” (ou plus précisément, tout foncteur “logiquement accessible”) entre catégories accessibles devrait être accessible. Ce n'est malheureusement pas vrai, puisqu'un composé de deux foncteurs Hom
\[
\Ens \longrightarrow \Ens\ , \qquad\quad E \longmapsto \underbrace{\Hom(\Hom(E,I),J)}_{J^{(I^E)}}
\]
\printedpage{180}
% head: A. Grothendieck : Dérivateurs, Ch. XVIII, Édition des Universités de Montpellier II et Paris VII
% foot: Version 17 décembre 2025
est non accessible, si $I$, $J$ sont tous deux de cardinal $\geqslant 2$. Prenant p.ex.\ $I = J = \{0,1\}$, on trouve le foncteur
\[
E \longmapsto \mathfrak{P}(\mathfrak{P}(E))\ ,
\]
qui n'est pas accessible.

\textbf{4.16.1.} Deuxième désillusion : j'avais cru que pour toute catégories accessible $\mathcal{M}$, toute % [sic] toute catégories -> toute catégorie, singular required after "toute"
sous-catégorie $\mathcal{N}$ strictement pleine et stable par facteurs directs “raisonnable” (p.ex. “logiquement accessible”) devrait être accessible dans $\mathcal{M}$. Il n'en est rien, même dans le cas particulier de l'image inverse de la sous-catégorie accessible dans Ens
\msp{465}
\[
\underbrace{\Ens^*}_{=\,\{E\,\in\Ens\,\mid\, E\neq\emptyset\}} \quad\subset\ \Ens
\]
par un contrafoncteur accessible \textit{représentable}
\[
\begin{aligned}
\mathcal{M} \quad &\longrightarrow \quad (\Ens)^\circ \\
x \quad &\longmapsto \quad \Hom_{\mathcal{M}}(x,z)\ ,
\end{aligned}
\]
i.e.\ la sous-catégorie ouverte (donc strictement pleine) $\mathcal{M}_z$ de $\mathcal{M}$ formée des objets $x$ tels que $\Hom_{\mathcal{M}}(x,z) \neq \emptyset$, n'est par toujours accessible dans $\mathcal{M}$. Ici, je ne suis pas parvenu à % [sic] n'est par toujours -> n'est pas toujours, misprint "par" for "pas"
un exemple aussi élémentaire (cf.\ appendice) -- j'ai pris $\mathcal{M} = (k\text{-alg.\ comm.})$, $k$ un corps infini et $z$ l'objet initial $k$ de $\mathcal{M}$ -- j'ai trouvé que $\mathcal{M}_z$ n'est pas $\mathrm{L}_\pi$-stable dans $\mathcal{M}$, quelque soit le cardinal $\pi$ % [sic] quelque soit -> quel que soit, standard spelling of the locution
-- mais cet exemple utilisait que $\mathfrak{U} = \mathfrak{U}_1$, plus petit univers contenant $\mathbf{N}$. J'ignore si pour tout autre univers, on peut trouver un corps $k$ de cardinal $\pi \in \mathfrak{U}$ assez grand, de façon que tout $\pi$-ultrafiltre sur un ensemble $E \in \mathfrak{U}$ soit trivial. Peut-être
\msp{466}
un tel exemple, où non la construction du couple $(\mathcal{M},\mathcal{N})$, mais la démonstration du fait que $\mathcal{N}$ n'est pas accessible dans $\mathcal{M}$, fait appel à des hypothèses particulières sur l'univers ambiant, ne doit pas être considéré comme (tout à fait !) “logiquement accessible”. Pour être vraiment probant, il faudrait un exemple de nature plus élémentaire.

% margin: a vertical bar is printed in the left margin beside the three lines from "L_pi-stable dans M pour pi assez grand" to "logiquement accessible) serait accessible dans M"
Il reste la possibilité qu'il y ait un théorème métamathématique qui assure que tout couple $(\mathcal{M},\mathcal{N})$ logiquement accessible, avec $\mathcal{M}$ accessible, et $\mathcal{N}$ sous-catégorie strictement pleine $\mathrm{L}_\pi$-stable dans $\mathcal{M}$ pour $\pi$ assez grand, soit accessible dans $\mathcal{M}$ (\footnote[188]{C'est faux, même pour une sous-catégorie fermée, cf.\ 4.16.7.10 ci-dessous.}). Ceci s'appliquerait en particulier au cas d'une sous-catégorie \textit{fermée} : toute sous-catégorie fermée (qui soit logiquement accessible) serait accessible dans $\mathcal{M}$. L'exemple-test le plus simple serait celui du complémentaire fermé $\mathcal{M}'_z$ de la sous-catégorie ouverte $\mathcal{M}_z$ ci-dessus (\footnote[189]{\textbf{N.B.} elles est accessible dans $\mathcal{M}$ si et seulement si $\mathcal{M}_z$ l'est ! % [sic] elles est -> elle est, agreement with singular verb
}), i.e.\ la catégorie des $x \in \Ob\mathcal{M}$ tels que $\Hom_{\mathcal{M}}(x,z) = \emptyset$, ou encore
\msp{467}
l'image inverse de la sous-catégorie fermée accessible
\[
e \;=\; \{\emptyset\} \;\subset\; (\Ens)
\]
réduite à $\{\emptyset\}$, par le foncteur représentable $x \longmapsto \Hom(x,z) : \mathcal{M} \longrightarrow \Ens$. Il se pourrait donc qu'une telle catégorie soit toujours accessible dans $\mathcal{M}$. Mais j'ai été incapable de le
\printedpage{181}
% head: A. Grothendieck : Dérivateurs, Ch. XVIII, Édition des Universités de Montpellier II et Paris VII
% foot: Version 17 décembre 2025
prouver, même dans le cas particulier ci-dessus, où $\mathcal{M} = (k\text{-alg.\ comm.})$, $k$ étant un corps commutatif infini. La difficulté est (dans tous les cas) la suivante : trouver, pour $z$ fixé, un cardinal $\pi > \pi_0$ tel que pour tout $x \in \Ob\mathcal{M}$, et tout système (ou \textit{un} système) $\pi_0$-$\pi$-adapté $(x_i)_{i\in I}$ de limite inductive $x_0$ ($\pi_0$ un cardinal donné, utile pour $\mathcal{M}$), \textit{si} $\Hom_{\mathcal{M}}(x,z)$ ($\simeq \varprojlim \Hom_{\mathcal{M}}(x_i,z)$) est vide, alors $\exists\ i_0 \in I$ tel que $\Hom_{\mathcal{M}}(x_{i_0},z)$ soit déjà vide. Dans l'exemple $\mathcal{M} = (k\text{-alg.\ comm.})$, il faut donc prouver que l'on peut trouver un cardinal $\pi$
tel que pour tout $k$-algèbre $A$ avec $\Hom_{\mathcal{M}}(A,k) = \emptyset$, il existe une sous-$k$-algèbre $A_i$ de % [sic] pour tout k-algèbre -> pour toute k-algèbre, "algèbre" is feminine
cardinal $\leqslant \pi$, telle que $\Hom_k(A_i,k) = \emptyset$ (\footnote[190]{C'est faux, cf.\ annotation marginale p.~466.}). J'ai tendance à croire que ça doit être vrai (et même sans
\msp{468}
supposer que $k$ soit un corps, et aussi en prenant la condition $\Hom_{\mathcal{M}}(A,Z) = \emptyset$, où $Z$ est un objet quelconque de $\mathcal{M}$). Ce genre de questions semble étonnement délicat \dots % [sic] étonnement -> étonnamment, adverb required

\textbf{4.16.2.} Il est possible, plus généralement, que la question suivante ait une réponse affirmative : Soit
\[
f \;:\; \mathcal{M} \longrightarrow \mathcal{N}^\circ
\]
un foncteur accessible, avec $\mathcal{M}$, $\mathcal{N}$ catégories accessibles, et soit $\underline{T}$ une sous-catégorie de $\mathcal{N}$ accessible dans $\mathcal{N}$, et telle qu'il existe un cardinal $\pi$ ayant les propriétés suivantes : $\mathcal{N}^\circ$ satisfait $\mathrm{L}_\pi$ (i.e.\ $\mathcal{N}$ stable par $\varprojlim$ grandes devant $\pi_1$). Alors $\underline{S} \defeq f^{-1}(\underline{T})$ est-elle accessible dans $\mathcal{M}$ (\footnote[191]{Faux, avec $\mathcal{N} = \Ens$ et un foncteur représentable $\mathcal{M} \longrightarrow \Ens^\circ$, prenant $I = \{\emptyset\}$, cf.\ annotations marginales p.~466, 467.}) ? Ceci généralise le cas $\mathcal{N} = (\Ens)$, $\underline{T} = \{\emptyset\}$, $f$ foncteur représenté par $z \in \Ob\mathcal{M}$, envisagé tout à l'heure. Avec le même type d'exemple, mais en prenant $\mathcal{N} = (\Ens)^*$
\msp{469}
au lieu de $\mathcal{N} = \{\emptyset\}$, on voit que la condition supplémentaire de stabilité sur $\underline{T}$ n'est pas
superflue. (\textbf{N.B.} Elle sont à assurer que $\underline{S}$ est $\mathrm{L}_\pi$-stable dans $\mathcal{M}$ pour $\pi$ grand.) % [sic] Elle sont -> Elles sont, agreement with plural verb

\textbf{4.16.3.} Questions apparentée : j'avais espéré que pour un foncteur accessible entre catégories accessibles % [sic] Questions apparentée -> Question apparentée, agreement
\[
f \;:\; \mathcal{N} \longrightarrow \mathcal{M}\ ,
\]
la $\natural$-image essentielle, formée des $y \in \Ob\mathcal{M}$ qui sont isomorphes à un facteur direct d'un objet de la forme $f(x)$, avec $x \in \Ob\mathcal{N}$, serait accessible dans $\mathcal{N}$. L'exemple ci-dessus (nullement élémentaire !) de la sous-catégorie $\mathcal{M}_z$ de $\mathcal{M}$ (formée des $x \in \Ob\mathcal{M}$ tels que $\Hom_{\mathcal{M}}(x,z) \neq \emptyset$) montre qu'il n'en est rien. En effet, $\mathcal{M}_z$ est l'image essentielle, et aussi l'image $\natural$-essentielle, du foncteur
\[
\mathcal{M}/z \longrightarrow \mathcal{M}\ ,
\]
or $\mathcal{M}/z$ est accessible, et le foncteur précédent accessible, si $\mathcal{M}$ est accessible (vérification par AQT \texttt{[âne qui trotte]}, cf.\ plus bas, énoncé plus général).

\msp{470}

On a l'impression que la raison pour laquelle on a de tels contre-exemples, c'est que l'on a des $\natural$-images essentielles qui ne sont $\mathrm{L}_\pi$-stables dans $\mathcal{M}$ pour aucun cardinal $\pi \in \mathfrak{U}$. Cela amène à modifier la conjecture ainsi : soit $\pi$ un cardinal tel que $\mathcal{M}$ satisfasse à $\mathrm{L}_\pi$,
\printedpage{182}
% head: A. Grothendieck : Dérivateurs, Ch. XVIII, Édition des Universités de Montpellier II et Paris VII
% foot: Version 17 décembre 2025
et soit $\underline{S}(\pi)$ la plus petite sous-catégorie strictement pleine de $\mathcal{M}$, $\mathrm{L}_\pi$-stable dans $\mathcal{M}$, et qui contienne l'image essentielle de $f$. (On pourrait appeler $\underline{S}(\pi)$ la $\mathrm{L}_\pi$-image essentielle.) Alors $\underline{S}(\pi)$ est accessible dans $\mathcal{M}$. De plus, la famille (décroissante) des $\underline{S}(\pi)$, pour $\pi$ variable, est stationnaire, i.e.\ il existe un cardinal $\pi$ tel que $\pi' \geqslant \pi$ implique $\underline{S}(\pi') = \underline{S}(\pi)$. Il est clair alors que $\underline{S}(\pi)$ est la plus petite sous-catégorie de $\mathcal{M}$ qui soit accessible dans $\mathcal{M}$ et qui contienne l'image essentielle de $f$. La même question se pose pour un foncteur $g^\circ : \mathcal{N}^\circ \longrightarrow \mathcal{M}$, où $g : \mathcal{N} \longrightarrow \mathcal{M}^\circ$ est un foncteur accessible, $\mathcal{N}$ et $\mathcal{M}$ accessibles. Je conjecture que la réponse est encore affirmative.

\msp{471}

% margin: a hand-drawn curved stroke is printed in the left margin beside the last three lines of the following paragraph ("pleine M de C^ qui soit ..." to "sans pourtant être accessible.")
Soit $\mathcal{M}$ une catégorie satisfaisant la condition L, et admettant une petite sous-catégorie pleine strictement génératrice $\mathcal{C}$, formée d'objets accessibles. On a donc un foncteur pleinement fidèle et accessible
\[
\mathcal{M} \;\into\; \mathcal{C}^\wedge\ ,
\]
$\mathcal{M}$ $\mathrm{L}_\pi$-stable dans $\mathcal{C}^\wedge$ pour $\pi$ convenable. Si $\overline{\mathcal{M}}$ est l'image essentielle, de sorte que $\mathcal{M} \simeq \overline{\mathcal{M}}$, c'est donc une sous-catégorie strictement pleine $\mathrm{L}_\pi$-stable dans $\mathcal{C}^\wedge$. Donc on s'attend à ce qu'elle soit accessible dans $\mathcal{C}^\wedge$ (si la situation décrite est “logiquement accessible”), ou ce qui revient au même, que $\mathcal{M}$ soit accessible. (Comparer 4.8.1, p.~187.) Mais il ne doit pas être difficile (\footnote[192]{OK en prenant un \textit{fermé} $\underline{S}$ non accessible dans $\mathcal{M}$ d'une catégorie accessible $\mathcal{M}$, qu'on plonge dans un $\mathcal{C}^\wedge$. Mais ces $\underline{S}$ n'a \textit{pas} une petite famille strictement génératrice ! % [sic] ces S n'a -> ce S n'a, plural determiner with singular verb
}), en procédant à un nombre $\pi_{\mathfrak{U}}$ de choix (où $\pi_{\mathfrak{U}} = \card\mathfrak{U}$), de faire quand même des exemples (“logiquement inaccessibles”) d'une sous-catégorie strictement pleine $\mathcal{M}$ de $\mathcal{C}^\wedge$ qui soit $\mathrm{L}_\pi$-stable, mais non accessible dans $\mathcal{C}^\wedge$. Cette $\mathcal{M}$ admet donc une petite sous-catégorie strictement génératrice formée d'objets accessibles (\footnote[193]{Ce n'est pas clair, il faudrait pour cela que le foncteur d'inclusion soit un adjoint à gauche !}) et satisfait à $\mathrm{L}_\pi$, sans pourtant être accessible.

\msp{472}

\textbf{4.16.4} (\footnote[194]{descendre après 4.16.5.}). Soit $\mathcal{M}$ une catégorie telle que $\mathcal{M}$ et $\mathcal{M}^\circ$ soient accessibles. Je conjecture qu'alors $\mathcal{M}$ est essentiellement petite. Plus généralement, supposons $\mathcal{M}$ accessible, et que $\mathcal{M}$ admette une petite sous-catégorie qui soit \textit{strictement co}génératrice. Alors $\mathcal{M}$ est-elle essentiellement petite ? Plus généralement, si $\mathcal{M}$ admet une petite sous-catégorie strictement génératrice, et aussi une petite sous-catégorie strictement cogénératrice, i.e.\ si $\mathcal{M}$ et $\mathcal{M}^\circ$ sont toutes deux équivalentes à des sous-catégories pleines d'une $\mathcal{C}^\wedge$, avec $\mathcal{C}$ petite, alors $\mathcal{M}$ est-elle essentiellement petite ?

\textbf{4.16.5.} Soit $I$ une petite catégorie, $I_\bullet$ le cône droit sur $I$ (obtenu en rajoutant un objet final $e$), $\mathcal{M}$ une catégorie accessible,
\[
\underline{\Hom}_{\square}(I_\bullet,\mathcal{M}) \;\into\; \underline{\Hom}(I_\bullet,\mathcal{M})
\]
la sous-catégorie strictement pleine de la catégorie (accessible) $\underline{\Hom}(I_\bullet,\mathcal{M})$, formée des foncteurs $f$ qui font de $f(e)$ la limite inductive dans $\mathcal{M}$ des $f(i)$ ($i \in \Ob I$). Cette sous-catégorie est $\mathrm{L}_\pi$-stable pour $\pi$ grand, donc il y a lieu
\msp{473}
de conjecturer qu'elle est accessible dans $\underline{\Hom}(I_\bullet,\mathcal{M})$. $\Big($ Je rappelle qu'on a vu dans 4.14 que c'est OK si les $\varinjlim$ de type $I$ sont représentables dans $\mathcal{M}$. Et qu'une réponse affirmative, dans le cas où
\printedpage{183}
% head: A. Grothendieck : Dérivateurs, Ch. XVIII, Édition des Universités de Montpellier II et Paris VII
% foot: Version 17 décembre 2025
\[
\begin{tikzcd}[row sep=small, column sep=small]
 & 1 \\
I = 0 \arrow[ur] \arrow[dr] & \\
 & 2
\end{tikzcd}
\quad ,\qquad\quad \text{donc} \qquad\quad
\begin{tikzcd}[row sep=small, column sep=small]
 & 1 \arrow[dr] & \\
I = 0 \arrow[ur] \arrow[dr] & & 3 \\
 & 2 \arrow[ur] &
\end{tikzcd}
\quad ,
\]
implique que la sous-catégorie $\underline{\text{Ép}}(\mathcal{M})$ dans $\underline{\mathrm{Fl}}(\mathcal{M})$ est accessible dans $\underline{\mathrm{Fl}}(\mathcal{M})$.$\Big)$

\textbf{N.B.} Si $\mathcal{M}$ admet une petite sous-catégorie \textit{co}génératrice, alors la question est justiciable de la conjecture 4.16.2 relative à un foncteur accessible $\mathcal{M} \longrightarrow \mathcal{N}^\circ$, ici $\underline{\Hom}(I_\bullet,\mathcal{M}) \longrightarrow \underline{\Hom}(I^\circ \times \mathcal{C},\Ens)^\circ \dots$

\textbf{4.16.6.} Il reste la question de l'accessibilité relative d'une catégorie $\mathcal{E}$ sur une petite catégorie $\mathcal{C}$. Je pense à présent, grâce à 4.15.9, et son corollaire 4.15.10, que l'énoncé suivant doit “tomber” par AQT \texttt{[âne qui trotte]} :

\textbf{Proposition 4.16.6.1.} \textit{Soit $f : \mathcal{E} \longrightarrow \mathcal{C}$ un foncteur, avec $\mathcal{C}$ petite catégorie karoubienne. On suppose que $f$ satisfait la condition $H' = H$ (1+2+3) de 4.15.1 (page 415), où il suffit H 2, H 3, pour toute flèche $\lambda$} non identique \textit{de $\mathcal{C}$. On suppose de plus que pour une catégorie $\mathcal{C}'$ d'un des deux types qu'on va décrire, et un foncteur $\mathcal{C}' \longrightarrow \mathcal{C}$, $\mathcal{E}' = \mathcal{E} \times_{\mathcal{C}} \mathcal{C}'$ soit accessible, et $\mathcal{E}' \longrightarrow \mathcal{C}'$ accessible.}

\msp{474}

$\Big($ \textit{On prend pour $\mathcal{C}'$ la catégorie}
\[
\mathcal{Q} = \left\{ \quad
\begin{tikzcd}[column sep=2.2em]
\xi_0 = \xi^\natural \arrow[r, shift left, "\sigma"] \arrow[dr, "u_0"'] & \xi \arrow[l, shift left, "\rho"] \arrow[d, "u = u_0\rho"] \arrow[loop right, "p"] \\
 & \eta
\end{tikzcd}
\qquad
\begin{array}{rclll}
\rho\sigma & = & \mathrm{id}_{\xi_0} & & \\[2ex]
\sigma\rho & = & p\ , & \text{\textit{d'où}} & \left\{\begin{array}{rcl} p^2 & = & p \\ p\sigma & = & \sigma \\ \rho p & = & \rho \end{array}\right. \\[4ex]
u & = & u_0\rho\ , & \text{\textit{d'où}} & \left\{\begin{array}{rcl} up & = & u \\ u\sigma & = & u_0 \end{array}\right.
\end{array}
\right.
\]
\textit{(trois objets $\xi$, $\xi_0$, $\eta$ et cinq flèches non identiques $\sigma$, $\rho$, $p$, $u_0$, $u$, liés par les relations ci-dessus), ou la catégorie (isomorphe à $Q^\circ$)} % Doubt resolved against PDF leaf 25: Q is ordinary mathematical italic, retained.
\[
\mathcal{Q}' = \left\{ \quad
\begin{tikzcd}
\xi_0 \arrow[r, shift left, "\sigma"] & \xi \arrow[l, shift left, "\rho"] \arrow[loop right, "p"] \\
 & \eta \arrow[ul, "u_0"] \arrow[u, "u"]
\end{tikzcd}
\qquad
\begin{array}{lcl}
\left.\begin{array}{rcl} \rho\sigma & = & \mathrm{id}_{\xi_0} \\ \sigma\rho & = & p \end{array}\right\} & \Longrightarrow & \left\{\begin{array}{rcl} p^2 & = & p \\ p\sigma & = & \sigma \\ \rho p & = & \rho \end{array}\right. \\[4ex]
\begin{array}{rcl} u & = & \sigma u_0 \end{array} & \Longrightarrow & \left\{\begin{array}{rcl} pu & = & u \\ \rho u & = & u_0 \end{array}\right.
\end{array}
\right.
\]
\textit{On se limite aux foncteurs $\mathcal{Q} \longrightarrow \mathcal{C}$, $\mathcal{Q}' \longrightarrow \mathcal{C}$ tels que $p$ soit transformé en une flèche} non identique \textit{de $\mathcal{C}$ (\footnote[195]{Je ne suis pas sûr si c'est raisonnable -- à vérifier avec soin !}) (donc la condition relative à $\mathcal{Q}$, $\mathcal{Q}'$ devient vide, si $\mathcal{C}$ est sans projecteurs).} $\Big)$

\textit{Sous ces conditions, $\mathcal{E}$ est accessible, et $f : \mathcal{E} \longrightarrow \mathcal{C}$ est accessible.}

% ---- from der-ch18_p186-190.tex
\setcounter{footnote}{195}

\printedpage{184}
% head: A. Grothendieck : Dérivateurs, Ch. XVIII, Édition des Universités de Montpellier II et Paris VII
% foot: Version 17 décembre 2025

\msp{page 475}
On trouve alors le

\begingroup\itshape
\noindent\textbf{\upshape Corollaire 4.16.6.2.} Soit $f : \cE \longrightarrow \cC$ un foncteur. Conditions équivalentes :
\begin{enumerate}[label={}]
\item[a)] Pour toute catégorie $\cC'$ \textup{petite} et \textup{karoubienne}, et tout foncteur $\cC' \longrightarrow \cC$, $\cE' = \cE \times_{\cC} \cC'$ est accessible, et le foncteur $\cE' \longrightarrow \cC'$ est accessible.
\item[b)] Comme a), mais $\cC'$ étant une des trois catégories $\Delta_1$, $\cQ$, $\cQ'$ (où $\cQ$, $\cQ'$ sont définies dans 4.15.6.1 ci-dessus).
\item[b$'$)] \textup{(}\footnote{?, à vérifier avec soin !}\textup{)}. (Si $\cC$ est sans projecteurs.) Comme b), mais en se limitant au cas $\cC' = \boldsymbol{\Delta}_1$ (à l'exclusion de $\cQ$, $\cQ'$).
\end{enumerate}
\endgroup

\begingroup\itshape
\noindent\textbf{\upshape Corollaire 4.16.6.3.} Supposons $f : \cE \longrightarrow \cC$ fibrant ou cofibrant. Alors les conditions équivalentes précédentes impliquent que :
\begin{enumerate}[label={}]
\item[1$^\circ$)] Les $\cE_s$ sont accessibles ($s \in \Ob \cC$).
\item[2$^\circ$)] Pour toute flèche $\lambda : s \longrightarrow t$ dans $\cC$, le foncteur $\lambda^* : \cE_t \longrightarrow \cE_s$ (resp. $\cE_s \longrightarrow \cE_t$) est accessible.
\end{enumerate}
\noindent D'autre part, l'ensemble de ces conditions 1$^\circ$, 2$^\circ$) équivaut à la validité de b) ci-dessus, dans le cas $\cC' = \boldsymbol{\Delta}_1$, ou encore à la validité
\msp{page 476}
de a) ci-dessus, quand $\cC'$ est \textup{sans projecteurs}.
\endgroup

\textbf{Question 4.16.6.4.} \textit{Les conditions 1$^\circ$ et 2$^\circ$ précédentes impliquent-elles celles de 4.16.6.2 ? En d'autres termes, si la catégorie base $\cC$ est égale à $\cQ$ ou à $\cQ'$, alors ces conditions impliquent-elles que $\cE$ est accessible, et $f : \cE \longrightarrow \cC$ accessible ?}

Ça semble idiot de ne pas connaître la réponse !

\textbf{4.16.6.5.} J'ai envie d'appeler catégorie \textit{accessible sur $\cC$}, une catégorie $\cE$ sur $\cC$ satisfaisant les conditions équivalentes de 4.16.6.1 -- du moins dans le cas où on suppose $\cC$ petite. Pour que cette terminologie soit vraiment satisfaisant, il faudrait que l'on sache prouver que ces conditions impliquent ceci : pour toute catégorie accessible $\cM$ et tout foncteur accessible $\cM \longrightarrow \cC$, la catégorie $\cE \times_{\cC} \cM$ est accessible, et le foncteur $\cE \times_{\cC} \cM \longrightarrow \cM$ est accessible. % [sic] satisfaisant -> satisfaisante, accord avec « terminologie »

\msp{page 476bis}
L'idée, c'est qu'on appellerait (sous hypothèse de petitesse) \textit{catégorie accessible $\cE$} sur une autre $\cC$, une catégorie ayant la propriété qu'on vient de dire -- mais en supposant maintenant de plus que $\cC$ satisfasse à $\mathrm{L}_\pi$, j'imagine, pour tout foncteur accessible $\cM \longrightarrow \cC$, avec $\cM$ accessible, $\cE \times_{\cC} \cM$ est accessible et le foncteur canonique $\cE \times_{\cC} \cM \longrightarrow \cE$ est accessible.

\msp{page 477}
\subsection*{4.16.7. Retour sur les sous-catégories fermées d'une catégorie accessible (4.16.1).}

\textbf{Proposition 4.16.7.1.} \textit{Soient $\cM$ une catégorie accessible, $(\uS_i)_{i\in I}$ une famille de sous-catégories fermées accessibles dans $\cM$. On ne suppose que \textup{card}\,$I$ soit petit. Soit $\uS$ la sous-catégorie strictement pleine de $\cM$ telle que $\Ob \uS = \bigcup_{i\in I} \Ob \uS_i$. Alors pour que $\uS$ % [sic] On ne suppose que -> On ne suppose pas que, « pas » apparemment omis
\printedpage{185}
% head: A. Grothendieck : Dérivateurs, Ch. XVIII, Édition des Universités de Montpellier II et Paris VII
% foot: Version 17 décembre 2025
soit accessible dans $\cM$, il faut et il suffit qu'il existe une \textup{petite} partie $J \subset I$, telle que $\Ob \uS = \bigcup_{i\in J} \uS_i$. (En particulier, si $I$ est ordonnée, et $i \longmapsto \uS_i$ croissante et grande devant tout petit cardinal, il faut et il suffit qu'il existe un $i \in I$ tel que $\uS = \uS_i$, i.e. tel que $j \geqslant i$ implique $\uS_j = \uS_i$.)}

C'est évidemment suffisant, en vertu de 4.13\dots{} Prouvons inversement que si $\uS$ est accessible dans $\cM$, alors la condition précédente est satisfaite. Soit $\pi$ un cardinal $\geqslant \pi_0$ tel que $\uS$ soit $\pi$-accessible dans $\cM$, donc
\[
\uS \ \xleftarrow{\ \sim\ }\ \Ind_\pi \uS(\pi) \ , \qquad\qquad \text{où } \uS(\pi) = \uS \cap \cM_\pi.
\]

\msp{page 478}
Soit $\uS(\pi)_0$ une sous-catégorie réduite de $\uS(\pi)$, telle que $\uS(\pi)_0 \hookrightarrow \uS(\pi)$ soit une équivalence, de sorte que $\uS(\pi)_0$ est petite. Pour tout $x \in \Ob \uS(\pi)_0$, soit $i_x \in I$ tel que $x \in \Ob \uS_{i_x}$. Soit $J$ l'ensemble de tous les $i_x$, que est une petite partie de $I$, et soit $\uS_J$ la sous-catégorie strictement pleine de $\cM$ ayant pour ensemble d'objets $\bigcup_{j\in J} \Ob \uS_j$. Donc $\uS(\pi)_0$ est contenue dans $\uS_J$, donc aussi $\uS(\pi)$, puisque $\uS_J$ est strictement pleine : % [sic] que est -> qui est, pronom relatif sujet
\[
\uS(\pi) \ \hookrightarrow\ \uS_J \ .
\]
D'ailleurs $\uS_J$ est accessible dans $\cM$ (loc. cit.), et est fermée dans $\cM$ comme réunion de sous-catégories fermées, \textit{a fortiori elle satisfait à $L_\pi$}. Donc, comme $\uS(\pi) \subset \uS_J$, le foncteur essentiellement surjectif $\Ind_\pi(\uS(\pi)) \longrightarrow \uS$ se factorise par $\uS_J$. Donc $\uS_J \hookrightarrow \uS$ est essentiellement surjectif, et comme $\uS_J$ est strictement pleine, on en conclut $\uS_J = \uS$, q.e.d.

\msp{page 479}
\textbf{Proposition 4.16.7.2.} \textit{Soient $\cM$ une catégorie accessible, $x$ un objet de $\cM$, et ${}_x\cM$ la sous-catégorie fermée de $\cM$ engendrée par $x$, formée des $y \in \Ob \cM$ tels que $\Hom(x,y) \neq \emptyset$, i.e. l'image essentielle de $x\backslash\cM \longrightarrow \cM$. Alors ${}_x\cM$ est une sous-catégorie de $\cM$ accessible dans $\cM$.}

En effet, c'est l'image inverse de la sous-catégorie $\mathrm{Ens}^*$ accessible dans $\mathrm{Ens}$ par le foncteur représentable (donc accessible)
\[
h_x \ :\ y \longmapsto \Hom(x,y) \ :\ \cM \ \longrightarrow\ \mathrm{Ens} \ ,
\]
donc est accessible dans $\cM$ en vertu de 4.13\dots

\textbf{Corollaire 4.16.7.3.} \textit{Soit $\cM$ une catégorie accessible, $\uS$ une sous-catégorie fermée. Pour que $\uS$ soit accessible, il faut et il suffit qu'il existe une \textup{petite} partie $J$ de $\uS$, telle que $\Ob \uS = \bigcup_{x\in J} \Ob\,(\,{}_x\cM)$.}

En effet, on a
\[
\Ob \uS \ =\ \bigcup_{x\in\Ob\uS} \Ob\,(\,{}_x\cM) \ ,
\]
i.e. $\uS$ apparaît comme réunion de
\msp{page 480}
la famille (pas nécessairement petite) de sous-catégories fermées ${}_x\cM$ de $\cM$ (indexées par $x \in \Ob \uS$), accessibles dans $\cM$ en vertu de 4.16.7.2. Donc notre assertion est un cas particulier de 4.16.7.1.

\printedpage{186}
% head: A. Grothendieck : Dérivateurs, Ch. XVIII, Édition des Universités de Montpellier II et Paris VII
% foot: Version 17 décembre 2025

On peut aussi, pour toute partie petite $J$ de $\Ob \uS$, considérer la sous-catégorie fermée $\uS_J \subset \uS$ de $\cM$ engendrée par $J$. L'ensemble $I$ des parties $J$ en question, ordonné par inclusion, est grand devant tout cardinal $\pi \in \U$, et $\uS$ est réunion croissante des $\uS_J$. D'autre part, la proposition nous dit que $\uS$ est accessible dans $\cM$ si et seulement s'il existe un $J \in I$, tel que $\uS_J = \uS$, i.e. tel que $\uS_J = \uS_{J'}$ pour $J \subset J'$.

\begingroup\itshape
\noindent\textbf{\upshape Proposition 4.16.7.4.} Soit $\cM$ une catégorie accessible. Les conditions suivantes sont équivalentes :
\begin{enumerate}[label={}]
\item[a)] Il existe une sous-catégorie $\uS$ de $\cM$ fermée dans $\cM$, et non accessible dans $\cM$.
\item[b)] Soit $I$ l'ensemble des ordinaux $\alpha$ tels que $\card \alpha \in \U$, i.e. des
\msp{page 481}
ordinaux $\alpha$ qui sont dans $\U$. Donc $I = I_{\alpha_0}$, où $\alpha_0$ est le plus petit ordinal qui n'appartienne pas à $\U$, et $I$ est grand devant tout petit cardinal $\pi$, i.e. devant tout cardinal dans $\U$ \textup{(}\footnote{introduire $I = I_{\U}$ avant l'énoncé de la proposition.}\textup{)}. La condition envisagée est la suivante :

Il existe une suite transfini % [sic] transfini -> transfinie, accord avec « suite »
\[
\underline{x} \ =\ (x_\alpha)_{\alpha\in I}
\]
d'éléments de $\Ob \cM$, telle que l'on ait
\[
\Hom(x_\beta, x_\alpha) \ =\ \emptyset \qquad\qquad \textit{si } \beta < \alpha.
\]
De façon précise, pour une telle suite transfinie, soit $\uS$ ou ${}_{\underline{x}}\cM$ la sous-catégorie fermée de $\cM$ engendrée par la famille des $x_\alpha$, i.e. telle que $\Ob \uS = \bigcup_{\alpha\in I} \Ob\ {}_{x_\alpha}\cM = \bigcup_{\alpha\in I} x_\alpha\backslash\Ob \cM$. Alors ${}_{\underline{x}}\cM$ est une sous-catégorie fermée non accessible dans $\cM$, et toute sous-catégorie fermée non accessible dans $\cM$ s'obtient de cette façon.
\end{enumerate}
\endgroup

\textsc{Démonstration}. Considérons une suite $\underline{x} = (x_\alpha)_{\alpha\in I}$ comme dessus. Si $J$ est une petite partie de $I$, il existe un $\alpha \in I$ qui majore les $\beta \in J$, donc $\Hom(x_\beta, x_\alpha) = \emptyset$, pour $\beta \in J$,
\msp{page 482}
ce qui signifie que $x_\alpha$ n'appartient à aucune des $\uS_\beta = {}_{x_\beta}\cM$, pour $\beta \in J$. Donc $\uS$ n'est par réunion d'une petite sous-famille de la famille des ${}_{x_\alpha}\cM$, donc en vertu de 4.16.7.1, $\uS$ n'est par accessible dans $\cM$. % [sic] n'est par réunion -> n'est pas réunion ; n'est par accessible -> n'est pas accessible, « par » pour « pas » (2 fois)

Inversement, soit $\uS$ une telle sous-catégorie dans $\cM$. Considérons un bon ordre sur $\Ob \uS$ (objets $s_\alpha$), et construisons la suite des $x_\alpha$ ($\geqslant s_\alpha$) par récurrence transfinie (\footnote{\textbf{N.B.} On a $\card \uS \notin \U$, sinon, $\uS$ serait petite, donc accessible dans $\cM$, contrairement à l'hypothèse.}).

1) On choisit $x_0 =$ premier élément de $\Ob \uS$ (existe, car $\Ob \uS \neq \emptyset$, i.e. $\uS \neq \emptyset$, car sinon, $\uS$ serait accessible dans $\cM$).

2) Si les $x_\beta$ ($\beta < \alpha$) sont construits, on prend pour $x_\alpha$ le premier élément de $\uS$ majorant les $x_\beta$ et qui n'appartient pas à la sous-catégorie fermée de $\cM$ engendrée par les $x_\beta$ ($\beta < \alpha$). (Cette sous-catégorie fermée est $\subset \uS$, et $\neq \uS$, car $\uS$ n'est pas accessible dans $\cM$.) On a donc bien
\[
\Hom(x_\beta, x_\alpha) \ =\ \emptyset \qquad\qquad \text{si } \beta < \alpha.
\]

\printedpage{187}
% head: A. Grothendieck : Dérivateurs, Ch. XVIII, Édition des Universités de Montpellier II et Paris VII
% foot: Version 17 décembre 2025

\msp{page 483}
Ainsi, la récurrence transfinie se poursuit sur $I$ tout entier. On voit par récurrence transfinie qu'on a, pour tout $\alpha \in I$,
\[
s_\alpha \ \leqslant\ x_\alpha
\]
(card $x_\alpha > x_\beta \geqslant s_\beta$ pour $\beta < \alpha$, donc $x_\alpha > s_\beta$ pour tout $\beta < \alpha$, ce qui équivaut à $x_\alpha \geqslant s_\alpha$). De plus, on trouve aussi, par récurrence transfinie, que
\[
\beta \ <\ \alpha \qquad\quad \Longrightarrow \qquad\quad
\text{\parbox[c]{5.5cm}{$s_\beta$ appartient à la sous-catégorie fermée $\uS_\alpha$ de $\cM$ engendrée par les $x_{\beta'}$ avec $\beta' < \alpha$.}}
\quad ,
\tag*{$(*)$}
\]
puisque par construction $s_\alpha$ est le \textit{premier} des éléments de $\uS$ qui n'est \textit{pas} dans ladite sous-catégorie $\uS_\alpha$ de $\cM$. Cela montre que tout $s_\beta$ appartient à la sous-catégorie fermée ${}_{\underline{x}}\cM$ de $\cM$ engendrée par la famille de \textit{tous} les $x_\alpha$, donc on trouve que $\uS \subset {}_{\underline{x}}\cM \subset \uS$, donc $\uS = {}_{\underline{x}}\cM$, q.e.d.

\begingroup\itshape
\noindent\textbf{\upshape Proposition 4.16.7.5.} Pour un univers $\U$ donné (contenant $\mathcal{N}$ comme élément), les conditions suivantes sont équivalentes :
\begin{enumerate}[label={}]
\item[1)] Il existe une catégorie $\U$-accessible $\cM$ ayant une sous-catégorie fermée non accessible % [sic] 1) -> a), la démonstration renvoie à a) et b)
\msp{page 484}
dans $\cM$.
\item[b)] Il existe une catégorie $\cC \in \U$, et une suite transfinie d'éléments de
\[
\cC^\wedge \ \overset{\text{\upshape déf}}{=}\ \underline{\mathrm{Hom}}(\cC, \mathrm{Ens}_{\U}) \ ,
\]
$\underline{x} = (x_\alpha)_{\alpha\in I}$, indexée par l'ensemble de tous les ordinaux $\alpha \in \U$, et telle que
\[
\Hom(x_\alpha, x_\beta) \ =\ \emptyset \qquad\qquad \textit{si } \alpha < \beta.
\]
\end{enumerate}
\endgroup

\textsc{Démonstration}. On a b) $\Longrightarrow$ a) en prenant $\cM = \cC^\wedge$, en vertu de la proposition précédente. On voit de même a) $\Longrightarrow$ b), en plongeant $\cM$ dans une catégorie $\cC^\wedge$ ($\cC$ étant une petite sous-catégorie strictement génératrice de $\cM$).

Je n'ai pas été capable (\footnote{Mais si, finalement, cf. 4.16.7.10 plus bas.}) de construire une suite $(x_\alpha)_{\alpha\in I}$ comme exigée dans cette proposition, ni de prouver qu'une telle suite n'existe pas, même dans un cas particulier tel que $\cM = (k\text{-alg. comm.})$, $k$ un anneau commutatif. Je n'ai pas
\msp{page 485}
même prouvé qu'on ne puisse trouver une suite transfinie de \textit{corps} $k_\alpha$, ayant les propriétés dites.

A fortiori, je n'ai pu tirer au clair la question si une sous-catégorie $\underline{\mathrm{L}}$-stable d'une catégorie $\cM$ accessible est toujours accessible dans $\cM$ (\footnote{C'est faux, cf. plus bas.}). Il se pourrait qu'il n'y ait pas de tel cas. Il se pourrait aussi qu'il y en ait, même avec une sous-catégorie fermée (laquelle est stable par tout type de $\varinjlim$ représentable dans $\cM$, sauf la $\varinjlim$ sur une catégorie d'indices vide $\ldots$).

\printedpage{188}
% head: A. Grothendieck : Dérivateurs, Ch. XVIII, Édition des Universités de Montpellier II et Paris VII
% foot: Version 17 décembre 2025

Pour les sous-catégories \textit{ouvertes}, je rappelle que j'ai construit (dans l'appendice) un exemple d'une telle sous-catégorie, non accessible dans $\cM$. Au sujet des sous-catégories ouvertes, je signale ici l'énoncé (quasi tautologique) :

\msp{page 486}
\textbf{Proposition 4.16.7.6} (\footnote{Cette proposition et son corollaire sont contenus dans 4.16.7.8 ci-dessous, il n'y a pas lieu d'en faire un énoncé séparé.}). \textit{Soient $\cM$ une catégorie accessible, $\uS$ une sous-catégorie \textup{ouverte} de $\cM$. Pour que $\uS$ soit accessible dans $\cM$, il faut et il suffit qu'elle soit $\underline{L}$-stable dans $\cM$, i.e. qu'il existe un cardinal $\pi_0$ tel que $\cM$ satisfasse à $L_{\pi_0}$, et que $\uS$ soit stable dans $\cM$ par $\varinjlim$ de type $L_{\pi_0}$. De façon plus précise, si cette condition est satisfaite, alors pour tout cardinal $\pi$ tel que $\pi \geqslant \pi_0$ et que $\pi$ soit utile pour $\cM$, $\uS$ est $\pi$-accessible dans $\cM$. Ou encore : si $\pi$ est un cardinal utile pour $\cM$, alors $\uS$ est $\pi$-accessible dans $\cM$ si et seulement s'il est $L_\pi$-stable dans $\cM$.}

Cette dernière assertion résulte aussitôt du corollaire suivant, lui-même tautologique.

\textbf{Corollaire 4.16.7.7.} \textit{Soit $\pi$ un cardinal comme indiqué ci-dessus, soit $\uS(\pi) = \uS \cap \cM_\pi$, et soit $x \in \Ob \cM$ et
\[
x \ \simeq\ \varinjlim_{i\in I} x_i \qquad\qquad
\begin{cases}
x_i \in \Ob \cM_\pi \\[1ex]
I \ \textit{grande devant}\ \pi
\end{cases}
\]
une représentation de $x$ comme image essentielle d'un objet de $\Ind_\pi \cM_\pi$. Alors $x \in \Ob \uS$ $\Longleftrightarrow$ $x_i \in \Ob \uS(\pi)$ $\forall\, i \in I$.}

\msp{page 487}
J'ai finalement tiré au clair la question de l'accessibilité des sous-catégories ouvertes et des sous-catégories fermées d'une catégorie accessible :

\begingroup\itshape
\noindent\textbf{\upshape Proposition 4.16.7.8.} Soient $\cM$ une catégorie accessible, $\uS$ une sous-catégorie fermée de $\cM$, $\uT$ la sous-catégorie ouverte complémentaire de $\uS$, $\pi$ un cardinal infini utile pour $\cM$. Les conditions suivantes sont équivalentes :
\begin{enumerate}[label={}]
\item[a)] $\uS$ est $\pi$-accessible dans $\cM$.
\item[a$'$)] $\uS$ est le fermé engendré par $\uS(\pi) \overset{\text{\upshape déf}}{=} \uS \cap \cM_\pi$, i.e. pour $x \in \Ob \cM$, on a $x \in \Ob \uS$ si et seulement si $\exists\, \xi \in \Ob \uS(\pi)$ tel que $\Hom_{\cM}(\xi, x) \neq \emptyset$.
\item[b)] $\uT$ est $\pi$-accessible dans $\cM$.
\item[b$'$)] $\uT$ est $\underline{L}_\pi$-stable dans $\cM$.
\end{enumerate}
\endgroup

\noindent Notons tout de suite le

\begingroup\itshape
\noindent\textbf{\upshape Corollaire 4.16.7.9.} Pour $\cM$, $\uS$, $\uT$ comme dessus, conditions équivalentes :
\begin{enumerate}[label={}]
\item[a)] $\uS$ est accessible dans $\cM$.
\item[a$'$)] Il existe une \textup{petite} partie $A$ de $\Ob \cM$, telle que $\uS$ soit la sous-catégorie fermée ${}_A\cM$ engendrée par $A$, i.e. $\uS = \bigcup_{\xi\in A} {}_\xi\cM$, où ${}_\xi\cM$ est la sous-catégorie fermée de $\cM$ formée des $x \in \Ob \cM$ tel que $\Hom(\xi, x) \neq \emptyset$.
\item[b)] \msp{page 488}$\uT$ est accessible dans $\cM$.
\end{enumerate}
\endgroup

% ---- from der-ch18_p191-195.tex
\printedpage{189}
% head: A. Grothendieck : Dérivateurs, Ch. XVIII, Édition des Universités de Montpellier II et Paris VII

\begin{enumerate}
\item[\textit{b$'$)}] \textit{$\underline{T}$ est $\underline{L}$-stable dans $\cM$.}
\end{enumerate}

\noindent\textbf{Corollaire 4.16.7.10.} \textit{Soit $\cM$ une catégorie accessible. Les conditions suivantes sont équivalentes} (\footnote[202]{Exemple où ces conditions sont satisfaites : $\cM = (k\text{-alg. comm.})$, $k$ un corps commutatif infini, $\mathfrak{U}$ étant le plus petit univers tel que $\mathbf{N}\in\mathfrak{U}$.}).
\begin{enumerate}
\item[\textit{a)}] \textit{Il existe une sous-catégorie fermée de $\cM$, non accessible dans $\cM$.}
\item[\textit{b)}] \textit{Il existe une sous-catégorie ouverte dans $\cM$, non accessible dans $\cM$.}
\item[\textit{c)}] \textit{Il existe une suite transfinie}
\[
\underline{x} \ = \ (x_\alpha)_{\alpha\in I}
\]
\textit{d'objets de $\cM$, indexée par l'ensemble $I$ des ordinaux de l'univers de référence $\mathfrak{U}$, telle que l'on ait}
\[
\Hom(x_\alpha,x_\beta) \ = \ \emptyset \qquad \text{\textit{pour} } \alpha<\beta .
\]
\end{enumerate}

\noindent En effet, l'équivalence de a) et b) est claire par le corollaire précédent, et celle de a) et c) n'est autre que 4.16.7.4.

\noindent\textsc{Démonstration de} 4.16.7.8. Nous aurons à utiliser le

\noindent\textbf{Lemme 4.16.7.11.} \textit{Soient $\cM$ une catégorie accessible, $\pi$ un cardinal utile pour $\cM$, $\underline{S}$ une sous-catégorie pleine de $\cM$, $\underline{S}(\pi)\defeq\underline{S}\cap\cM_\pi$, d'où un foncteur pleinement fidèle}
\[
\varphi \ : \ \Ind_\pi(\underline{S}(\pi)) \longrightarrow \cM \ .
\]

\msp{page 489}
\begin{enumerate}
\item[\textit{a)}] \textit{Soit $x\in\Ob\cM$. Pour que $x$ appartienne à l'image essentielle du foncteur $\psi$, il faut et il suffit que la sous-catégorie pleine $\underline{S}(\pi)/x$ de la catégorie $\cM_\pi/x$ (qui est filtrante et même grande devant $\pi$) soit cofinale (donc aussi grande devant $\pi$).} % [sic] \psi -> \varphi, the functor of the lemma is called \varphi in the display above
\item[\textit{b)}] \textit{Pour que $\underline{S}$ soit $\pi$-accessible dans $\cM$, il faut et il suffit que pour tout $x\in\Ob\cM$, on ait}
\[
x \ \in \ \Ob\underline{S} \qquad\iff\qquad \underline{S}(\pi)/x \text{ \textit{est cofinale dans} } \cM_\pi/x
\]
\textit{(ou encore, que $\underline{S}$ satisfasse à $L_\pi$, et que pour tout $x$ dans $\cM$, on ait}
\[
x \ \in\Ob\underline{S} \qquad\Longrightarrow\qquad \underline{S}(\pi)/x \text{ \textit{cofinale dans} } \cM_\pi/x\text{\textit{.)}}
\]
\end{enumerate}

\noindent\textsc{Démonstration}. La partie b) est une conséquence tautologique de a), que je vais prouver. Dire que $x$ est dans l'image essentielle de $\underline{S}$, signifie que
\[
x \ \simeq \ \varinjlim_{i\in I} x_i \ , \tag*{$(*)$}
\]
avec $I$ grande devant $\pi$ et les $x_i$ dans $\underline{S}(\pi)$. On sait que $(*)$ implique ($I$ étant filtrante, et $\cM$ étant équivalente à $\Ind_\pi(\cM)$ puisque $\pi$ utile pour $\cM$) que $I\longrightarrow\cM_\pi/x$ est cofinale (4.4.2, page 79). Or ce foncteur se factorise par le foncteur d'inclusion pleinement fidèle $\underline{S}(\pi)/x\incl\cM_\pi/x$, lequel est cofinal si et seulement s'il satisfait la condition F 1. Il est immédiat que cette condition est satisfaite quand elle est satisfaite pour le foncteur composé en question, ce qui est bien le cas
\printedpage{190}
% head: A. Grothendieck : Dérivateurs, Ch. XVIII, Édition des Universités de Montpellier II et Paris VII
\msp{page 490}
par hypothèse.

\noindent Inversement, supposons que $\underline{S}(\pi)/x$ soit cofinale dans $\cM_\pi/x$. Comme on a
\[
x \xleftarrow{\ \sim\ } \varinjlim_{\cM_\pi/x} x_i \ ,
\]
$\pi$ étant utile pour $\cM$, on aura aussi par cofinalité
\[
x \xleftarrow{\ \sim\ } \varinjlim_{\underline{S}(\pi)/\pi} x_i \ , % [sic] \underline{S}(\pi)/\pi -> \underline{S}(\pi)/x, the cofinal subcategory is \underline{S}(\pi)/x
\]
or $\cM_\pi$ étant grande devant $\pi$, il en est de même de la sous-catégorie pleine cofinale $\underline{S}(\pi)/x$. Donc $x$ est dans l'image essentielle de $\Ind_\pi\underline{S}(\pi)$, q.e.d.

\noindent\textbf{Lemme 4.16.7.12.} \textit{Soient $I$ une catégorie filtrante, $J$ une sous-catégorie fermée, $K$ la sous-catégorie ouverte complémentaire. On a alors les équivalences}
\[
J \ \neq \ \emptyset \iff J \text{ \textit{cofinale} } \iff K \text{ \textit{non cofinale} } \iff K \ \neq \ I \ .
\]

\noindent\textsc{Démonstration} immédiate.

\noindent Je reviens à la démonstration de 4.16.7.8. On applique le lemme 4.16.7.11 à $\cM$, $\pi$, et aux deux
\msp{page 491}
sous-catégories $\underline{S}$, $\underline{T}$ de $\cM$, en notant que pour $x\in\Ob\cM$, $\underline{S}(\pi)/x$ est une sous-catégorie fermée de $\underline{S}(\pi)$, dont le complémentaire ouvert est $\underline{T}(\pi)/x$, de sorte qu'on peut appliquer le lemme 4.16.7.12. On trouve :
\[
\begin{aligned}
\underline{S}\ \pi\text{-accessible} \quad &\overset{4.16.7.11}{\iff} \quad \forall\ x\ \in\ \Ob\cM,\ \text{on a } (x\ \in\ \Ob\underline{S}\iff\underline{S}(\pi)/x \text{ cofinale dans}\\
&\qquad\quad \cM_\pi/x)\\
&\iff \quad \forall\, x\in\Ob\cM,\ \text{on a } \Bigl( \qquad \underbrace{x\notin\Ob\underline{S}}_{\text{(i.e. $x\in\Ob\underline{T}$, par 4.16.7.12)}}\\
&\qquad\qquad\qquad \iff \qquad \underbrace{\underline{S}(\pi)/x \text{ non cofinale dans } \cM_\pi/x}_{\text{(i.e. $\underline{T}(\pi)/x$ cofinale dans $\cM_\pi$, par 4.16.7.12)}} \ \Bigr)\\
&\overset{4.16.7.11}{\iff} \quad \underline{T} \text{ est } \pi\text{-accessible.}
\end{aligned}
\]
On a donc l'équivalence de a) et b). D'autre part, $\underline{S}(\pi)/x$ cofinale dans $\cM_\pi/x$ signifie simplement $\underline{S}(\pi)/x\neq\emptyset$ (4.16.7.12), donc $\underline{S}$ $\pi$-accessible signifie que $\underline{S}$ est formée des $x\in\Ob\cM$ tels que $\underline{S}(\pi)\neq\emptyset$, ce qui signifie aussi que $\underline{S}$ est le fermé dans $\cM$ engendré par $\underline{S}(\pi)$, d'où a $\iff$ a$'$. D'autre part, l'équivalence b $\iff$ b$'$ n'est autre que 4.6.7.6 (dernière formulation). % [sic] 4.6.7.6 -> 4.16.7.6, the surrounding references are all to 4.16.7.x

\msp{page 492}

\section*{4.17\quad Retour sur les paratopos, et sur les opérations $\cM\boxtimes\cN$, $\underline{\operatorname{Hom}}_{!}(\cM,\cN)$}

\noindent\texttt{[Ce paragraphe n'existe pas.]}

\printedpage{191}
% head: A. Grothendieck : Dérivateurs, Ch. XVIII, Édition des Universités de Montpellier II et Paris VII

\msp{page 1}

\section*{Appendice}

\subsection*{Ultraproduits de corps, et catégories accessibles}

\noindent Soit $(k_\xi)_{\xi\in E}$ une famille de corps, indexé par un ensemble $E$. $\Bigl($ Soit $\Phi$ un filtre non trivial dans $E$, i.e. une partie $\Phi\subset\mathfrak{P}(E)$ telle que :
\begin{enumerate}
\item[a)] $\emptyset\notin\Phi$, i.e. $A\in\Phi\Longrightarrow A\neq\emptyset$.
\item[b)] $A,B\in\Phi\Longrightarrow A\cap B\in\Phi$.
\item[c)] Si $A\in\Phi$, $A\subset B$, alors $B\in\Phi$.
\item[d)] $\Phi\neq\emptyset$ ($\iff E\in\Phi$).
\end{enumerate}
(Cela revient à dire que $\Phi$ est une partie fermée de l'ensemble ordonné $\mathfrak{P}(E)$, \textit{filtrante à gauche}, et ne contient pas le plus petit élément de $\mathfrak{P}(E)$, i.e. non égale à $\mathfrak{P}(E)$ tout entier.) $\Bigr)$ Pour toute partie $J$ de $E$, posons
\[
k_J \ = \ \prod_{\xi\in J} k_\xi \ ,
\]
on a donc un système projectif
\[
(k_J)_{J\in\mathfrak{P}(E)} \ , \qquad \text{i.e. } J\subset J' \text{ donne } k_{J'}\longrightarrow k_J \ ,
\]
et le système
\[
(k_J)_{J\in\Phi}
\]
est un système inductif d'annaux, indexé par l'ensemble ordonné filtrant $\Phi^o$. Posons % [sic] d'annaux -> d'anneaux, misspelling
\[
k_\Phi \ = \ \varinjlim_{J\in\Phi} k_J \ .
\]
Si $\Psi$ est un filtre plus fin que $\Phi$, i.e.
\msp{page 2}
tel que $\Psi\supset\Phi$, on trouve un homomorphisme canonique
\[
k_\Phi \longrightarrow k_\Psi \ ,
\]
donc si $\cF(E)=\cF$ est l'ensemble des filtres sur $E$, ordonnés par inclusion, on trouve un système inductif
\[
\Phi \longmapsto k_\Phi \ : \ \cF \longrightarrow \Ann \ .
\]
On n'a pas utilisé encore que les $k_\xi$ sont des corps, en fait il suffirait que ce soient des objets d'une catégorie $\cA$ stable par petits produis et par petits $\varinjlim$ filtrantes. Prenons maintenant $\cA=\Ann$ (catégorie des anneaux unitaires, non nécessairement commutatifs), et supposons que les $k_\xi$ soient des corps. % [sic] produis -> produits, misspelling

\printedpage{192}
% head: A. Grothendieck : Dérivateurs, Ch. XVIII, Édition des Universités de Montpellier II et Paris VII

\noindent\textbf{Proposition 1.} \textit{Supposons que les $k_\xi$ ($\xi\in E$) soient des} corps (\footnote[203]{NB Cette condition est nécessaire si on veut que $k_\Phi$ soit un corps pour tout ultrafiltre.}). \textit{Pour que $k_\Phi$ soit un corps, il faut et il suffit que $\Phi$ soit un ultrafiltre, i.e. que pour toute décomposition $E=A\amalg B$, on ait $A\in\Phi$ ou $B\in\Phi$.} (NB Si $A\in\Phi$, on a $B\notin\Phi$, puisque sinon, on aurait $A\cap B\in\Phi$, donc $\emptyset\in\Phi$, absurde.)

\noindent Nécessité : Soit, pour toute partie $A$ de $E$,
\[
f_A\in k_E \qquad \text{défini par} \qquad f_A(\xi) \ = \ \left\{
\begin{array}{lll}
1_{k_\xi} & \text{si} & \xi\in A\\[1ex]
0_{k_\xi} & \text{si} & \xi\in E\backslash A \ .
\end{array}\right.
\]

\msp{page 3}
\noindent Soit
\[
f_{A,\Phi} \ \in k_\Phi
\]
l'image de $f_A$ par l'homomorphisme canonique $k_E\longrightarrow k_\Phi$. Si $E=A\amalg B$, on aura
\[
f_A \ + \ f_B \ = \ 1 \ , \qquad f_A\, f_B \ = \ 0 \ ,
\]
d'où
\[
f_{A,\Phi} \ + \ f_{B,\Phi} \ = \ 1 \ , \qquad f_{A,\Phi}\, f_{B,\Phi} \ = \ 0 \ .
\]
Si $k_\Phi$ est un corps, cela implique qu'on a $f_{A,\Phi}=0$ ou $f_{B,\Phi}=0$. Mais si p.ex. $f_{A,\Phi}=0$, on voit que cela signifie qu'il existe un $J\in\Phi$ tel que $f_A|J=0$, ce qui signifie aussi que $J\subset B$. De même, si $f_{B,\Phi}=0$, on aura $J\subset A$. Donc $\Phi$ est un ultrafiltre.

\noindent Inversement, supposons que $\Phi$ soit un ultrafiltre. Prouvons que $k_\Phi$ est un corps, i.e. que pour $f\in k_\Phi$, $f\neq\emptyset\Longrightarrow f$ inversible. Notons que $f$ provient d'un élément d'un $k_J$ ($J\in\Phi$), donc est l'image canonique d'un $g\in k_E$. Dire que $f=0$ signifie que $\exists\ J\in\Phi$ tel que $g|J=0$. L'inverse $f\neq0$ signifie donc que pour tout $J\in\Phi$, on a $g|J\neq\emptyset$.

\msp{page 4}
\noindent Soit alors $A$ l'ensemble des $\xi\in E$ tel que $g(\xi)\neq\emptyset$, $B=E\backslash A$ l'ensemble des $\xi\in E$ tels que $g(\xi)\neq\emptyset$. On a donc $B\notin\Phi$ puisque $f\neq0$, et comme $\Phi$ est un ultrafiltre, on a donc $A\in\Phi$. Mais soit $g'\in k_E$ défini par $g'(\xi)=f(\xi)^{-1}$ si $\xi\in A$, $g'(\xi)=1$ si $\xi\in B$, et soit $f'$ l'image de $g'$ dans $k_\Phi$. On voit alors qu'on a % [sic] B = E\A ... tels que g(\xi) \neq \emptyset -> g(\xi) = 0, the complement of A is the set where g vanishes
\[
f\ f' \ = \ f' \ = \ f \ = \ 1 \ ,
\]
ce qui montre bien que $k_\Phi$ est un corps.

\noindent Supposons maintenant que tous les $k_\xi$ ($\xi\in E$) soient égaux à un même corps $k$. Alors les $k_J$ forment un système projectif de $k$-algèbres, les $k_\Phi$ un système inductif de $k$-algèbres. Si $k$ est commutatif, il en est de même des $k_J$, $k_\Phi$, et en particulier pour $\Phi$ un ultrafiltre, $k_\Phi$ est une \textit{extension} du corps $k$.

\noindent Revenons au cas général, mais les $k_\xi$ commutatifs. Soit $\cF_0\subset\cF$ l'ensemble des ultrafiltres sur $E$. Notons que pour tout $\Phi\in\cF$, $k_E\longrightarrow k_\Phi$ est surjectif. Donc si $\Phi\in\cF$, $k_\Phi$ s'identifie à un corps quotient de l'anneau $k_E$.

\printedpage{193}
% head: A. Grothendieck : Dérivateurs, Ch. XVIII, Édition des Universités de Montpellier II et Paris VII

\msp{page 5}
\noindent Pour tout anneau $\cA$, soit $\Max(\cA)$ le spectre maximal, i.e. l'ensemble des idéaux maximaux de $\cA$. On trouve donc une application canonique
\[
\left\{
\begin{array}{lll}
\cF_0 & \longrightarrow & \Max(k_E)\\[1ex]
\ \Phi & \longmapsto & \Ker(k_E\longrightarrow k_\Phi) \ \defeq \ \mathfrak{m}_\Phi \ .
\end{array}\right. \tag*{$(*)$}
\]

\noindent\textbf{Proposition 2.} \textit{L'application précédente est} bijective.

\noindent Prouvons d'abord qu'elle est injective, cela résulte aussitôt du

\noindent\textbf{Corollaire 1} (\footnote[204]{\textbf{NB} Pour $\Phi\in\cF$ quelconque, $\mathfrak{m}_\Phi\defeq\Ker(\chi_\Phi : k_E\longrightarrow k_\Phi)$ détermine $\Phi$, puisque
\[
A \ \in \ \Phi \qquad\iff\qquad \chi_\Phi(f_A) \ = \ 1
\]
(mais il faut supposer les $K_\xi\neq0$, i.e. $0_\xi\neq1_\xi\ \ldots$).}). \textit{Soit, pour $\Phi\in\cF_0$,}
\[
\chi_\Phi \ : \ k_E \longrightarrow k_\Phi
\]
\textit{l'application canonique. Alors pour $A\in\mathfrak{P}(E)$, on a}
\[
A \ \notin \ \Phi \qquad\iff\qquad \chi_\Phi(f_A) \ = \ 0 \ , \quad \text{\textit{i.e.} } f_A\in\mathfrak{m}_\Phi
\]
\textit{(i.e. $A\in\Phi\iff\chi_\Phi(f_A)\neq0$, i.e. $f_A\notin\mathfrak{m}_\Phi$)} (\footnote[205]{\textbf{NB} On a donc
\[
(\chi_\Phi(f_A) \ \neq \ 0) \qquad\iff\qquad (\chi_\Phi(f_A) \ = \ 1) \ .
\]}).

\noindent En effet, $\chi_\Phi(f_A)=0$ signifie qu'il existe $B\in\Phi$ tel que $f_A|B=0$, i.e. tel que $B\subset E\backslash A$ ; mais cela équivaut à $E\backslash A\in\Phi$, donc à $A\notin\Phi$ puisque $\Phi$ est un ultrafiltre, q.e.d.

\msp{page 6}
\noindent Reste à prouver que l'application $(*)$ est surjective. Soit donc $\mathfrak{m}$ un idéal maximal de $k_E$, et soit
\[
\Phi \ \defeq \ \{A\in\mathfrak{P}(E)\mid\chi(f_A) \ = \ 1\} \ ,
\]
où $\chi : k_E\longrightarrow k_E/\mathfrak{m}=\kappa(\mathfrak{m})$ est l'application canonique. Je dis que $\Phi$ est un ultrafiltre. En effet :
\begin{enumerate}
\item[a)] $\emptyset\notin\Phi$ car $f_\emptyset=0$, d'où $\chi(f_\emptyset)=0\neq1$ dans $k_\Phi$.
\item[b)] Si $A,B\in\Phi$, on a $A\cap B\in\Phi$, car $f_{A\cap B}=f_A\cdot f_B$, d'où $\chi(f_{A\cap B})=\chi(f_A)\chi(f_B)=1$.
\item[c)] Si $A\subset B$ et $A\in\Phi$, alors $B\in\Phi$. En effet, on a $f_A=f_Af_B$, d'où $\chi(f_A)=\chi(f_A)\chi(f_B)$, et comme $\chi(f_A)=1$, on trouve $\chi(f_B)=1$.
\item[d)] On a $E\in\Phi$, i.e. $\chi(f_E)=1$, or $f_E=1$, donc $\chi(f_E)=1$.
\end{enumerate}
Ainsi, $\Phi$ est un filtre -- on n'a utilisé que le fait que $A/\mathfrak{m}$ est un anneau $\neq0$ (pour conclure que $\emptyset\notin\Phi$). Prouvons que si $A/\mathfrak{m}$ est intègre, $\Phi$ est un \textit{ultrafiltre}, i.e. :
\begin{enumerate}
\item[e)] Si $A\notin\Phi$, alors $B\defeq E\backslash A\in\Phi$.
\end{enumerate}
En effet, on a $f_A+f_B=1$, $f_Af_B=0$, d'où

% ---- from der-ch18_p196-200.tex
\printedpage{194}
% head: A. Grothendieck : Dérivateurs, Ch. XVIII, Édition des Universités de Montpellier II et Paris VII

\noindent\hspace*{1em}\textcircled{a}\quad $\chi(f_A)+\chi(f_B)=1$, et\\
\hspace*{1em}\textcircled{b}\quad $\chi(f_A)\chi(f_B)=0$.\\
$\chi(f_A)\neq 1$, on conclut de \textcircled{a} $\chi(f_B)\neq 0$, d'où par \textcircled{b} $\chi(f_A)=0$, d'où de \textcircled{a} $\chi(f_B)=1$, d'où $B\in\Phi$.

\msp{page 7}
Prouvons enfin que $\mathfrak{m}=\mathfrak{m}_\Phi$. Comme $\mathfrak{m}_\Phi$ est un idéal maximal, il suffit de prouver
\[
\mathfrak{m}_\Phi\ \subset\ \mathfrak{m}\ .
\]
Or si $A\in\mathfrak{P}(E)$ est telle que $A\notin\Phi$, ou $f_A\in\mathfrak{m}_\Phi$ (cor.~1), on a aussi $f_A\in\mathfrak{m}$. Donc il suffit de prouver le

\textbf{Lemme.} \textit{Soit $\Phi\in\mathcal{F}_0$. Alors $\mathfrak{m}_\Phi$ est l'idéal engendré par les $f_A$ pour $A\in\mathfrak{P}(E)\backslash\Phi$. Plus précisément, les $f_A$ sont dans $\mathfrak{m}_\Phi$, et tout $f\in\mathfrak{m}_\Phi$ est de la forme $ff_A$, avec $A$ comme dessus.}

En effet, $f\in\mathfrak{m}_\Phi$ signifie qu'il existe $B\in\Phi$ tel que $f|B=0$. Mais cela signifie aussi que $f=ff_A$, où $A=E\backslash A$ \texttt{[plutôt} $A=E\backslash B$\texttt{]}, d'où $A\notin\Phi$.

Cela prouve la proposition 2, et en même temps :

\textbf{Corollaire 2.} \textit{Tout anneau quotient intègre de $k_E=\prod_{\xi\in E}k_\xi$ est un corps.}

Notons maintenant que si $\Phi\in\mathcal{F}$ est un filtre quelconque sur $E$, alors $k_\Phi$ est un anneau quotient de $k_E$, donc
\[
\operatorname{Max}(k_\Phi)\ \hookrightarrow\ \underbrace{\operatorname{Max}(k_E)}_{\simeq\,\mathcal{F}_0}\ . \tag*{$(*)$}
\]
Ceci posé, on trouve :

% a vertical bar is printed in the left margin beside the statement of Corollaire 3
\textbf{Corollaire 3} (\footnote[206]{Marche pour tout système d'anneaux $(k_\xi)_{\xi\in E}$ \textit{non nuls}, et $\Phi,\Psi\in\mathcal{F}$, $\Psi$ plus fin que $\Phi$, i.e. $\Phi\subset\Psi$ si et seulement si $\mathfrak{m}_\Phi\subset\mathfrak{m}_\Psi$ (pas la peine que $\Psi$ soit un ultrafiltre).}). \textit{L'inclusion précédente identifie} $\operatorname{Max}(k_\Phi)$ \textit{à l'ensemble des ultrafiltres $\Psi$ sur $E$ qui sont} plus fins \textit{que $\Phi$.}

En effet, si $\Psi$ est un tel ultrafiltre, on
\msp{page 8}
a le diagramme commutatif
\[
\begin{tikzcd}
k_E \arrow[r, "\chi_\Phi"] \arrow[rr, bend right=15, "\chi_\Psi"'] & k_\Phi \arrow[r] & k_\Psi
\end{tikzcd}\ , \tag*{$(**)$}
\]
ce qui montre que $\Psi$ \texttt{[plutôt} $\mathfrak{m}_\Psi$\texttt{]} est dans l'image de $(*)$. Inversement, supposons que $\Psi\in\mathcal{F}_0$ provienne de $\operatorname{Max}(k_\Phi)$, i.e. que l'on ait $\mathfrak{m}_\Phi\subset\mathfrak{m}_\Psi$ (où $\mathfrak{m}_\Phi$ est le noyau de $\chi_\Phi$), ou encore qu'on ait une factorisation $(**)$ de $\chi_\Psi$. Prouvons que $\Psi$ est plus fin que $\Phi$, i.e. que $\Psi\supset\Phi$. Soit donc $A\in\Phi$, prouvons que l'on a $A\in\Psi$. Mais $A\in\Phi$ implique \texttt{[le]} fait que $\chi_\Phi(f_A)=1$, donc son image dans $k_\Psi$ est égale à 1, donc on a $A\in\Psi$ par le corollaire 1.

\printedpage{195}
% head: A. Grothendieck : Dérivateurs, Ch. XVIII, Édition des Universités de Montpellier II et Paris VII

\msp{page 9}
Supposons à présent que tous les $k_\xi$ sont égaux \texttt{[à]} un même anneau \fbox{$\ k\neq 0\ $}\,. Donc $k_J$, $k_\Phi$ sont des `anneaux sous $k$' $k\longrightarrow k_J\longrightarrow k_\Phi$, d'ailleurs $k\longrightarrow k_J$ (si $J\neq\emptyset$) et $k\longrightarrow k_\Phi$ sont injectifs. On se pose la question quand $k\longrightarrow k_\Phi$ est \textit{bijectif}, i.e. surjectif. Comme $k_\Phi\simeq k_E/\mathfrak{m}_\Phi$, cela signifie donc que $k_E=k+\mathfrak{m}_\Phi$, i.e. que tout $f\in k_E$ s'écrit
\[
f\ =\ c\ +\ g\qquad
\begin{cases}
c\ \in\ k\ , & \text{i.e. } c(\xi)=\underbrace{c_0}_{\in\,k}\ \forall\,\xi\in E\ ,\\[2ex]
g\ \in\ \mathfrak{m}_\Phi & \text{i.e. } \exists\,A\in\Phi \text{ avec } g|A=0\ .
\end{cases}
\]
Mais pour $f\in k_E$, $c\in k$ et $A\in\Phi$ fixés, l'existence de $g\in k_E$ telle que $g|A=0$ et $f=c+g$ équivaut à $f|A=c_A$ (fonction constante dans $A$, de valeur $c$), i.e. $A\subset f^{-1}(\{c\})$. Pour $f$, $c$ fixés, l'existence de $A\in\Phi$ satisfaisant la relation équivaut à $f^{-1}(\{c\})\in\Phi$. Et pour $f$ fixé, l'existence de $c$ équivaut à celle d'une fibre $E_c$ de $f:E\longrightarrow k$, qui soit dans $\Phi$. Ainsi :

\textbf{Proposition 3.} \textit{Soit $\Phi\in\mathcal{F}$. Pour que $k\longrightarrow k_\Phi$ soit surjectif, il faut et il suffit que pour toute famille de parties} essentiellement disjointes \textit{$(E_\xi)_{\xi\in k}$ de $E$ indexée par $k$, de réunion $E$, il existe un $\xi\in k$ tel que $E_\xi\in\Phi$. A fortiori, il} faut \textit{que $\Phi$ soit un ultrafiltre (puisque} $\operatorname{card}k\geqslant 2$ \textit{par hypothèse).}

\msp{page 10}
En effet, la donnée d'une telle famille $(E_\xi)_{\xi\in k}$ équivaut à la donnée d'une application $E\longrightarrow k$, i.e. d'un élément de $f_E$. % [sic] $f_E$ -> $k_E$, an application $E\to k$ is an element of $k_E$

\textbf{Scholie.} La propriété qu'on vient d'expliciter pour un filtre $\Phi$ sur $E$, ne dépend de $k$ que par l'intermédiaire de $\pi=\operatorname{card}k$. On peut aussi l'énoncer en disant que pour toute \textit{partition} $(E_i)_{i\in I}$ de $E$ en un ensemble $I$ des parties tel que $\operatorname{card}I\leqslant\pi$, il existe un (et un seul) des $E_i$ qui soit dans $\Phi$. On dit alors que $\Phi$ est un $\pi$-\textit{ultrafiltre.}

Si $\pi\geqslant 2$ est fini, cette propriété équivaut à dire que $\Phi$ est un ultrafiltre. Par ailleurs, quel que soit $\pi$, les ultrafiltres $\Phi$ \textit{triviaux} (i.e. tels qu'il existe $\xi\in E$ avec $\{\xi\}\in\Phi$, ou encore tel que $\Phi=\{A\in\mathfrak{P}(E)\mid\xi\in A\}$) sont des $\pi$-ultrafiltres. Si $\pi$ est infini, il n'est pas clair qu'il existe des $\pi$-ultrafiltres non triviaux. On sait qu'il n'en existe sur $E$ \textit{que} si $\operatorname{card}E\notin\mathfrak{U}_1$, où $\mathfrak{U}_1$ est le plus petit univers tel que $\mathbf{N}\in\mathfrak{U}_1$, i.e. tel
\msp{page 11}
que $\mathfrak{U}$ contienne comme élément un ensemble infini.

Si $\mathfrak{U}=\mathfrak{U}_1$, et si $k$ est un corps infini, p.ex. $k=\mathbf{Q}$, alors pour tout ensemble $E\in\mathfrak{U}$ et pour tout ultrafiltre \textit{non trivial} $\Phi$ sur $E$, $k_E$ est une extension \textit{non triviale} de $k$. Mais si $\mathfrak{U}=\mathfrak{U}_\alpha$ est un univers quelconque, et $\pi$ le premier cardinal dans $\mathfrak{U}$ qui n'appartienne pas à un univers $\mathfrak{U}'\subset\mathfrak{U}$ strictement plus petit, je présume que pour tout $E\in\mathfrak{U}$, tout $\pi$-ultrafiltre sur $E$ est trivial. % [sic] $k_E$ est une extension -> $k_\Phi$ est une extension, the extension of $k$ attached to the ultrafilter $\Phi$ is $k_\Phi$

\printedpage{196}
% head: A. Grothendieck : Dérivateurs, Ch. XVIII, Édition des Universités de Montpellier II et Paris VII

\subsection*{Application à un exemple d'une sous-catégorie non accessible.}

Soit $\mathcal{M}$ une catégorie $\mathfrak{U}$-accessible, et $a\in\operatorname{Ob}\mathcal{M}$, d'où un foncteur
\[
\begin{aligned}
f^o\ :\ \mathcal{M}^o\ &\longrightarrow\ \mathrm{Ens}\\
x\ &\longmapsto\ \operatorname{Hom}_{\mathcal{M}}(x,a)\ \overset{\text{déf}}{=}\ a(x)\ ,
\end{aligned}
\]
foncteur commutant aux $\limproj$ quelconques, donc tel que $f:\mathcal{M}\longrightarrow(\mathrm{Ens})^o$ commute aux $\limind$ quelconques, et en particulier $f$ est accessible. Considérons la sous-catégorie strictement pleine
\msp{page 12}
\[
\underbrace{\mathrm{Ens}^*}_{=\,\mathrm{Ens}\backslash\{\emptyset\}}\ \subset\ \mathrm{Ens}
\]
des ensembles non vides, on sait qu'elle est accessible, mais elle n'est par stable par $\limproj$ de type $\mathrm{L}_\pi$, quelque grand qu'on prenne $\pi$. % [sic] n'est par stable -> n'est pas stable, typo

Considérons
\[
\underbrace{\underline{S}\ =\ \underline{S}_a}_{\overset{\text{déf}}{=}\,f^{-1}(\mathrm{Ens}^*)}\ \subset\ \mathcal{M}\ ,
\]
sous-catégorie strictement pleine de $\mathcal{M}$ formée des $x\in\operatorname{Ob}\mathcal{M}$ tels que
\[
a(x)\ \neq\ \emptyset\ ,\qquad\quad \text{i.e. } x\leqslant a \text{ pour le préordre canonique de } \mathcal{M}\ .
\]
C'est donc une sous-catégorie \textit{ouverte} de $\mathcal{M}$. Son complémentaire fermé $\underline{S}'=\underline{S}'_a$ est la sous-catégorie fermée image inverse, par $f$, de la sous-catégorie strictement pleine de Ens réduite au seul objet $\{\emptyset\}$. Cette sous-catégorie de Ens est la catégorie finale type. Elle est stable dans Ens par $\limind$ quelconques, et par $\limproj$ filtrantes quelconques. Ainsi, $\underline{S}'$ est stable dans $\mathcal{M}$ par $\limind$ filtrantes quelconques. Ainsi, $\underline{S}'$ est stable dans $\mathcal{M}$ par $\limind$ filtrantes quelconques, en particulier, elle est $\mathrm{L}_\pi$-stable pour $\pi$ assez grand. Explicitons :

\msp{page 13}
\[
\begin{array}{cl}
\underbrace{\underline{S}'\ =\ \underline{S}'_a}_{}\ \subset\ \mathcal{M} & \\
\Big\|\ \text{\scriptsize déf} & \text{sous-catégorie fermée de } \mathcal{M}\\
f^{-1}(\{\emptyset\})\ = & \text{formée des } x\ \in\ \operatorname{Ob}\mathcal{M} \text{ tels}\\
 & \text{que } a(x)=\emptyset.
\end{array}
\]

Comme $\underline{S}'$, étant fermée dans $\mathcal{M}$, est stable par $\limind$ non vides et en particulier $\mathrm{L}_\pi$-stable pour $\pi$ grand, il y a plus de raison de s'attendre à ce que $\underline{S}'$ soit accessible dans $\mathcal{M}$ que pour son complémentaire $\underline{S}$, dont il n'est pas clair qu'il ait cette propriété de stabilité, mais seulement stabilité par $\limproj$ non vides, et en particulier par facteurs directs.

En fait, nous allons voir qu'il peut arriver que $\underline{S}$ et $\underline{S}'$ soient tous deux non accessibles dans $\mathcal{M}$ :

\textbf{Proposition 4.} \textit{Soient $\mathfrak{U}$ un univers, contenant un cardinal infini $\pi_0$ tel que pour tout $E\in\mathfrak{U}$, tout $\pi_0$-ultrafiltre sur $E$ soit trivial. Soit alors $k\in\mathfrak{U}$ un corps de cardinal $\pi'_0\geqslant\pi_0$, et soit $\mathcal{M}$ la catégorie}
\printedpage{197}
% head: A. Grothendieck : Dérivateurs, Ch. XVIII, Édition des Universités de Montpellier II et Paris VII
\msp{page 14}
\textit{des $k$-algèbres commutatives. Donc $\mathcal{M}$ est un $\mathfrak{U}$-paratopos, admettant $k$ comme objet initial. Soit $\underline{S}'$ la sous-catégorie strictement pleine, stable par petites $\limind$ non vides, formée des objets $A$ de $\mathcal{M}$ tels que $\operatorname{Hom}_{\mathcal{M}}(A,k)=\emptyset$, et $\underline{S}$ la sous-catégorie strictement pleine complémentaire, formée des objets $A$ tels que $\operatorname{Hom}_{\mathcal{M}}(A,k)=\emptyset$. Alors $\underline{S}'$ est une sous-catégorie} fermée \textit{dans $\mathcal{M}$, et $\underline{S}$ la sous-catégorie ouverte complémentaire, donc $\underline{S}$ (tout comme $\underline{S}'$) est stable par facteurs directs. Ceci posé :} % [sic] second $\operatorname{Hom}_{\mathcal{M}}(A,k)=\emptyset$ -> $\operatorname{Hom}_{\mathcal{M}}(A,k)\neq\emptyset$, $\underline{S}$ is the complement of $\underline{S}'$
\begin{enumerate}[label={}]
\item[\textit{a)}] \textit{Il n'existe pas de cardinal $\pi$ tel que $\underline{S}$ soit $L_\pi$-stable dans $\mathcal{M}$, a fortiori $\underline{S}$ n'est pas accessible dans $\mathcal{M}$.}
\item[\textit{b)}] (\footnote[207]{?, pas prouvé.}). \textit{Bien que $\underline{S}'$ soit $L_0$-stable dans $\mathcal{M}$ (i.e. stable par $\limind$ filtrantes, et même, on l'a dit, par toute $\limind$ non vide), $\underline{S}'$ n'est pourtant pas accessible dans $\mathcal{M}$.}
\end{enumerate}
\noindent\textsc{Démonstration}. L'assertion a) va résulter du

\textbf{Corollaire 1.} \textit{Sous les conditions de la proposition 4, soient $E\in\mathfrak{U}$ un ensemble infini, $\Phi_0$ le filtre des complémentaires des parties finies de $E$. Alors $\operatorname{Hom}_{\mathcal{M}}(k_{\Phi_0},k)=\emptyset$, a fortiori, pour tout filtre $\Phi$ sur $E$ plus fin que $\Phi_0$, on a $\operatorname{Hom}_{\mathcal{M}}(k_\Phi,k)=\emptyset$.}

\msp{page 15}
On a posé ici, bien sûr, pour tout filtre $\Phi$ sur $E$,
\[
k_\Phi\ =\ \limind_{J\in\Phi}k_J\ .
\]
On sait que pour toute $k$-algèbre commutative $\mathcal{A}$, $\operatorname{Hom}(\mathcal{A},k)$ est en correspondance 1-1 avec la partie de $\operatorname{Max}(\mathcal{A})$ formée des idéaux maximaux $\mathfrak{m}$ de $\mathcal{A}$ tels que $k\longrightarrow\mathcal{A}/\mathfrak{m}$ soit un isomorphisme. Si $\mathcal{A}$ est de la forme $k_\Phi$, $\Phi$ un filtre sur $E$, cet ensemble est en correspondance biunivoque avec l'ensemble des $\pi'_0$-ultrafiltres $\Psi$ sur $E$ plus fins que $\Phi$, où $\pi'_0=\operatorname{card}(k)$ (proposition 3, et cor.~3 de prop.~2, page 7). Comme $\pi'_0\geqslant\pi_0$, un tel ultrafiltre est trivial par hypothèse. Donc on trouve
\[
\operatorname{Hom}_{\mathcal{M}}(k_\Phi,k)\ \simeq\ \text{ensemble des } \xi\in E \text{ tels que } \{\xi\}=\textstyle\bigcap_{A\in\Phi}A\ . \tag*{$(*)$}
\]
Mais dire que $\Phi$ est plus fin que $\Phi_0$ signifie aussi que pour tout $\xi\in E$, on a $E\backslash\{\xi\}\in\Phi$, i.e. $\exists\,A\in\Phi$ avec $A\subset E\backslash\{\xi\}$, i.e. $\{\xi\}\notin\bigcap_{A\in\Phi}A$. Donc on a alors $\operatorname{Hom}_{\mathcal{M}}(k_\Phi,k)=\emptyset$.

La conclusion du corollaire 1 s'exprime par la relation
\[
k_\Phi\ \notin\ \operatorname{Ob}\underline{S}\ ,\qquad\quad \text{i.e. } k_\Phi\in\operatorname{Ob}\underline{S}'\ ,
\]
d'autre part, on a par définition
\[
k_\Phi\ =\ \limind_{J\in\Phi}k_J\ ,
\]
\msp{page 16}
et il est clair (comme ces $J\in\Phi$ sont $\neq\emptyset$), que les $k_J$ en question sont dans $\underline{S}$. Ainsi, pour un filtre $\Phi$ comme dans le corollaire 1, on trouve un système inductif d'objets de $\underline{S}$, indexé par $\Phi^o$, dont la limite inductive n'est pas dans $\underline{S}$. Donc le a) de la proposition 4 se déduit en notant que pour tout cardinal $\pi\in\mathfrak{U}$, on peut trouver un ensemble $E\in\mathfrak{U}$, et un filtre $\Phi$ sur $E$ plus fin que $\Phi_0$, tel que l'ensemble ordonné $\Phi^o$ soit grand devant $\pi$. Pour trouver un tel $\Phi$, il suffit de prendre $E$ de cardinal $\pi'>\pi$, et pour $\Phi$ l'ensemble
\printedpage{198}
% head: A. Grothendieck : Dérivateurs, Ch. XVIII, Édition des Universités de Montpellier II et Paris VII
des complémentaires des parties de $E$ de cardinal $\leqslant\pi$, de sorte que $\Phi^o$ est isomorphe à l'ensemble ordonné des parties de $E$ de cardinal $\leqslant\pi$, lequel est bien grand devant $\pi$. Cela prouve le a) de la proposition 4.

Prouvons maintenant le b). Notons le

\textbf{Lemme.} \textit{Soit $\mathcal{M}$ une catégorie accessible stable par $\limind$ de type $I_\alpha$, $\alpha$ un ordinal limite donné, $\underline{S}'$ une sous-catégorie fermée, $\pi$ un cardinal utile}
\msp{page 17}
\textit{pour $\mathcal{M}$, tel que} $\operatorname{card}\alpha\leqslant\pi$\textit{. Pour tout $x\in\operatorname{Ob}\underline{S}'$, choisissons une représentation $\omega$-$\pi$-adaptée ($\omega$ étant le premier ordinal infini)}
\[
x\ \simeq\ \limind_{i\in I_x}x_i
\]
\textit{(donc $I_x$ grande devant $\pi$, stable par $\limind$ de type $I_\alpha$, $i\longmapsto x_i:I\longrightarrow\mathcal{M}$ commute auxdites limites, enfin les $x_i\in\operatorname{Ob}\mathcal{M}_\pi$). Pour que $\underline{S}'$ soit $\pi$-accessible dans $\mathcal{M}$, il faut et il suffit que pour tout $x\in\operatorname{Ob}\underline{S}'$, il existe un $i\in I_x$ tel que $x_i\in\operatorname{Ob}\underline{S}'$.}

En effet, notons que $\underline{S}'$ étant fermée, est stable dans $\mathcal{M}$ par toute $\limind$ non vide représentable dans $\mathcal{M}$, et en particulier par $\limind$ de type $I_\alpha$. Donc on peut appliquer 4.13.22 (p.~342) à $\mathcal{M}$, $\underline{S}'$, $\pi$ et à $\Phi=\{I_\omega\}$. Alors a) $\Longrightarrow$ c) implique que si $\underline{S}'$ est $\pi$-accessible dans $\mathcal{M}$, alors pour tout $x\in\operatorname{Ob}\underline{S}'$, l'ensemble $I_{x,S'}$ des $i\in I_x$ tels que $x_i\in\operatorname{Ob}\underline{S}$ est cofinal dans $I_x$. Cela implique qu'il est non vide. (Et réciproquement d'ailleurs, $\underline{S}'$ étant un fermé dans $\mathcal{M}$, puisque si $x_i\in\underline{S}'$, alors pour $j\in\operatorname{Ob}I_x$ avec $i\leqslant j$, on a $x_j\in\underline{S}'$, i.e. $I_{x,\underline{S}'}$ contient l'image essentielle de $i\backslash I_x$, laquelle est cofinale dans $I$ \dots) Cela prouve le `il faut'. Prouvons le `il suffit'. Donc prouvons que $\underline{S}'$ est égale à l'image essentielle de $\operatorname{Ind}_\pi(\underline{S}'(\pi))$. Elle la contient, puisque $\underline{S}'$ est stable par $\limind$ de type $\mathrm{L}_\pi$. % Doubt resolved against PDF leaf 40: the first S-prime subscript is not underlined, retained. % [sic] $x_i\in\operatorname{Ob}\underline{S}$ -> $x_i\in\operatorname{Ob}\underline{S}'$, the set $I_{x,\underline{S}'}$ concerns the $x_i$ in $\underline{S}'$

\msp{page 18}
Donc il reste à prouver que tout $x\in\operatorname{Ob}\underline{S}'$ est dans l'image essentielle. Or l'hypothèse nous assure que $x\simeq\limind\limits_{i\in I_x}x_i$ est isomorphe à $\limind\limits_{i\in I_{x,\underline{S}'}}x_i$ (car $I_{x,\underline{S}'}$ est cofinale dans $I_x$), où cette fois les $x_i$ sont dans $\underline{S}'$. Comme $I_x$ grande devant $\pi$, de même la sous-catégorie pleine cofinale $I_{x,\underline{S}'}$, donc on gagne.

\textbf{Corollaire.} \textit{Sous les conditions précédentes, si $\underline{S}'$ est $\pi$-accessible, il en est de même de l'ouvert complémentaire $\underline{S}$,} pourvu \textit{que $\underline{S}$ soit stable dans $\mathcal{M}$ par $\limind$ de type $L_\pi$.}

Il faut prouver que $\underline{S}$ est égale à l'image essentielle de $\operatorname{Ind}_\pi\underline{S}(\pi)$. L'hypothèse `pourvu' assure que cette image essentielle est $\subset\underline{S}$. Inversement, tout [\,?] $x\in\operatorname{Ob}\underline{S}$ est dans cette image essentielle. On aura
\[
x\ \simeq\ \limind_{i\in I}x_i\ ,\qquad\quad x_i\in\operatorname{Ob}\mathcal{M}_\pi,\ I=\mathcal{M}_\pi/x,\ \text{grande devant } \pi\ ,
\]
et comme $x\in\operatorname{Ob}\underline{S}$ et $\underline{S}$ ouverte, les $x_i$ sont dans $\underline{S}$, donc dans $\underline{S}(\pi)$, OK.

Dans le cas qui nous occupe, où $\mathcal{M}=(k\text{-alg.})$, comme $\underline{S}$ n'est justement \textit{pas} stable dans $\mathcal{M}$ par $\mathrm{L}_\pi$ (quel que soit $\pi$) (c'est cela qui nous faisait conclure que
\msp{page 19}
$\underline{S}$ n'est pas accessible dans $\mathcal{M}$), on ne peut donc conclure de cette inaccessibilité de $\underline{S}$, celle de $\underline{S}'$. Notons le

% ---- from der-ch18_p201-201.tex
\printedpage{199}
% head: A. Grothendieck : Dérivateurs, Ch. XVIII, Édition des Universités de Montpellier II et Paris VII
% foot: Version 17 décembre 2025

\noindent\textbf{Corollaire.} \textit{Soit $\pi$ un cardinal infini. Pour que la sous-catégorie fermée $\underline{S}'$ de $\mathcal{M}$ (formée de $k$-algèbres commutatives $A$ telles que $\Hom_{\mathcal{M}}(A,k)=\emptyset$) soit $\pi$-accessible, il faut et il suffit que pour tout $k$-algèbre $A$ (de cardinal ${}>\pi$), telle que $\Hom_{\mathcal{M}}(A,k)=\emptyset$, existe une sous-algèbre $A_i$} de cardinal $\leqslant\pi$, \textit{telle que $\Hom_{\mathcal{M}}(A_i,k)=\emptyset$.} % [sic] tout k-algèbre -> toute k-algèbre, « algèbre » est féminin

\noindent(On applique le lemme, en prenant pour $\alpha$ le premier ordinal infini.)

\noindent\textbf{Question.} \textit{Cela est-il vérifié, en prenant simplement $\pi\geqslant\card k$~?} La question de l'existence d'un cardinal $\pi$, ayant la propriété indiquée plus haut, me paraît délicate, et j'ignore la réponse. Si on veut construire un contre-exemple pour tout $\pi$, les développements précédents suggèrent de prendre pour $A$ un produit infini d'extensions non triviales de $k$, ou un anneau $k_{\Phi}$ comme dans la démonstration de la proposition 4 a).
\end{document}
